% Môn: Vật lý
% Lớp: 10
% Chương: Chương V. Động lượng
% Bài: L10C5 Bài 29. Định luật bảo toàn động lượng
% Số câu: 162
% App đọc file này trực tiếp — sửa rồi Commit trên GitHub, bấm Đồng bộ GitHub .tex

% L10C5 Bài 29. Định luật bảo toàn động lượng

% ===== Câu 1 =====
% ID: AIQ_L10C5_FULL_0163
% Mức: NB
\dangbt{Nhận biết hệ kín và điều kiện bảo toàn động lượng}
\begin{ex}
% ID: AIQ_L10C5_FULL_0163
% Mức: NB
Một hệ kín ban đầu đứng yên tách thành hai phần. Phần thứ nhất có động lượng $6\ \mathrm{kg\,m/s}$ theo chiều dương. Độ lớn động lượng phần thứ hai bằng bao nhiêu?
\choice
{\True $6\ \mathrm{kg\,m/s}$}
{$5\ \mathrm{kg\,m/s}$}
{$3\ \mathrm{kg\,m/s}$}
{$12\ \mathrm{kg\,m/s}$}
\loigiai{
Hệ kín có tổng động lượng ban đầu bằng 0 nên $\vec p_2=-\vec p_1$. Suy ra kết quả $6\ \mathrm{kg\,m/s}$.
}
\end{ex}

% ===== Câu 2 =====
% ID: AIQ_L10C5_FULL_0164
% Mức: NB
\begin{ex}
% ID: AIQ_L10C5_FULL_0164
% Mức: NB
Một hệ kín ban đầu đứng yên tách thành hai phần. Phần thứ nhất có động lượng $12\ \mathrm{kg\,m/s}$ theo chiều dương. Độ lớn động lượng phần thứ hai bằng bao nhiêu?
\choice
{$7\ \mathrm{kg\,m/s}$}
{\True $12\ \mathrm{kg\,m/s}$}
{$6\ \mathrm{kg\,m/s}$}
{$24\ \mathrm{kg\,m/s}$}
\loigiai{
Hệ kín có tổng động lượng ban đầu bằng 0 nên $\vec p_2=-\vec p_1$. Suy ra kết quả $12\ \mathrm{kg\,m/s}$.
}
\end{ex}

% ===== Câu 3 =====
% ID: AIQ_L10C5_FULL_0165
% Mức: TH
\begin{ex}
% ID: AIQ_L10C5_FULL_0165
% Mức: TH
Một hệ kín ban đầu đứng yên tách thành hai phần. Phần thứ nhất có động lượng $20\ \mathrm{kg\,m/s}$ theo chiều dương. Độ lớn động lượng phần thứ hai bằng bao nhiêu?
\choice
{$9\ \mathrm{kg\,m/s}$}
{$10\ \mathrm{kg\,m/s}$}
{\True $20\ \mathrm{kg\,m/s}$}
{$40\ \mathrm{kg\,m/s}$}
\loigiai{
Hệ kín có tổng động lượng ban đầu bằng 0 nên $\vec p_2=-\vec p_1$. Suy ra kết quả $20\ \mathrm{kg\,m/s}$.
}
\end{ex}

% ===== Câu 4 =====
% ID: AIQ_L10C5_FULL_0166
% Mức: TH
\begin{ex}
% ID: AIQ_L10C5_FULL_0166
% Mức: TH
Một hệ kín ban đầu đứng yên tách thành hai phần. Phần thứ nhất có động lượng $30\ \mathrm{kg\,m/s}$ theo chiều dương. Độ lớn động lượng phần thứ hai bằng bao nhiêu?
\choice
{$11\ \mathrm{kg\,m/s}$}
{$15\ \mathrm{kg\,m/s}$}
{$60\ \mathrm{kg\,m/s}$}
{\True $30\ \mathrm{kg\,m/s}$}
\loigiai{
Hệ kín có tổng động lượng ban đầu bằng 0 nên $\vec p_2=-\vec p_1$. Suy ra kết quả $30\ \mathrm{kg\,m/s}$.
}
\end{ex}

% ===== Câu 5 =====
% ID: AIQ_L10C5_FULL_0167
% Mức: VD
\begin{ex}
% ID: AIQ_L10C5_FULL_0167
% Mức: VD
Một hệ kín ban đầu đứng yên tách thành hai phần. Phần thứ nhất có động lượng $14\ \mathrm{kg\,m/s}$ theo chiều dương. Độ lớn động lượng phần thứ hai bằng bao nhiêu?
\choice
{\True $14\ \mathrm{kg\,m/s}$}
{$9\ \mathrm{kg\,m/s}$}
{$7\ \mathrm{kg\,m/s}$}
{$28\ \mathrm{kg\,m/s}$}
\loigiai{
Hệ kín có tổng động lượng ban đầu bằng 0 nên $\vec p_2=-\vec p_1$. Suy ra kết quả $14\ \mathrm{kg\,m/s}$.
}
\end{ex}

% ===== Câu 6 =====
% ID: AIQ_L10C5_FULL_0168
% Mức: VD
\begin{ex}
% ID: AIQ_L10C5_FULL_0168
% Mức: VD
Một hệ kín ban đầu đứng yên tách thành hai phần. Phần thứ nhất có động lượng $24\ \mathrm{kg\,m/s}$ theo chiều dương. Độ lớn động lượng phần thứ hai bằng bao nhiêu?
\choice
{$11\ \mathrm{kg\,m/s}$}
{\True $24\ \mathrm{kg\,m/s}$}
{$12\ \mathrm{kg\,m/s}$}
{$48\ \mathrm{kg\,m/s}$}
\loigiai{
Hệ kín có tổng động lượng ban đầu bằng 0 nên $\vec p_2=-\vec p_1$. Suy ra kết quả $24\ \mathrm{kg\,m/s}$.
}
\end{ex}

% ===== Câu 7 =====
% ID: AIQ_L10C5_FULL_0169
% Mức: VD
\begin{ex}
% ID: AIQ_L10C5_FULL_0169
% Mức: VD
Một hệ kín ban đầu đứng yên tách thành hai phần. Phần thứ nhất có động lượng $36\ \mathrm{kg\,m/s}$ theo chiều dương. Độ lớn động lượng phần thứ hai bằng bao nhiêu?
\choice
{$13\ \mathrm{kg\,m/s}$}
{$18\ \mathrm{kg\,m/s}$}
{\True $36\ \mathrm{kg\,m/s}$}
{$72\ \mathrm{kg\,m/s}$}
\loigiai{
Hệ kín có tổng động lượng ban đầu bằng 0 nên $\vec p_2=-\vec p_1$. Suy ra kết quả $36\ \mathrm{kg\,m/s}$.
}
\end{ex}

% ===== Câu 8 =====
% ID: AIQ_L10C5_FULL_0170
% Mức: VDC
\begin{ex}
% ID: AIQ_L10C5_FULL_0170
% Mức: VDC
Một hệ kín ban đầu đứng yên tách thành hai phần. Phần thứ nhất có động lượng $50\ \mathrm{kg\,m/s}$ theo chiều dương. Độ lớn động lượng phần thứ hai bằng bao nhiêu?
\choice
{$15\ \mathrm{kg\,m/s}$}
{$25\ \mathrm{kg\,m/s}$}
{$100\ \mathrm{kg\,m/s}$}
{\True $50\ \mathrm{kg\,m/s}$}
\loigiai{
Hệ kín có tổng động lượng ban đầu bằng 0 nên $\vec p_2=-\vec p_1$. Suy ra kết quả $50\ \mathrm{kg\,m/s}$.
}
\end{ex}

% ===== Câu 9 =====
% ID: AIQ_L10C5_FULL_0171
% Mức: VDC
\begin{ex}
% ID: AIQ_L10C5_FULL_0171
% Mức: VDC
Một hệ kín ban đầu đứng yên tách thành hai phần. Phần thứ nhất có động lượng $22\ \mathrm{kg\,m/s}$ theo chiều dương. Độ lớn động lượng phần thứ hai bằng bao nhiêu?
\choice
{\True $22\ \mathrm{kg\,m/s}$}
{$13\ \mathrm{kg\,m/s}$}
{$11\ \mathrm{kg\,m/s}$}
{$44\ \mathrm{kg\,m/s}$}
\loigiai{
Hệ kín có tổng động lượng ban đầu bằng 0 nên $\vec p_2=-\vec p_1$. Suy ra kết quả $22\ \mathrm{kg\,m/s}$.
}
\end{ex}

% ===== Câu 10 =====
% ID: AIQ_L10C5_FULL_0172
% Mức: VDC
\begin{ex}
% ID: AIQ_L10C5_FULL_0172
% Mức: VDC
Một hệ kín ban đầu đứng yên tách thành hai phần. Phần thứ nhất có động lượng $36\ \mathrm{kg\,m/s}$ theo chiều dương. Độ lớn động lượng phần thứ hai bằng bao nhiêu?
\choice
{$15\ \mathrm{kg\,m/s}$}
{\True $36\ \mathrm{kg\,m/s}$}
{$18\ \mathrm{kg\,m/s}$}
{$72\ \mathrm{kg\,m/s}$}
\loigiai{
Hệ kín có tổng động lượng ban đầu bằng 0 nên $\vec p_2=-\vec p_1$. Suy ra kết quả $36\ \mathrm{kg\,m/s}$.
}
\end{ex}

% ===== Câu 11 =====
% ID: AIQ_L10C5_FULL_0173
% Mức: NB
\begin{ex}
% ID: AIQ_L10C5_FULL_0173
% Mức: NB
Xét các phát biểu thuộc dạng “Nhận biết hệ kín và điều kiện bảo toàn động lượng” ở tình huống số 1.
\choiceTF
{\True Hệ kín có tổng ngoại lực bằng 0 hoặc không đáng kể.}
{\True Trong hệ kín, tổng động lượng được bảo toàn.}
{Mọi hệ vật đều là hệ kín.}
{\True Nội lực không làm thay đổi tổng động lượng hệ.}
\loigiai{
\begin{itemchoice}
\item (Đúng) Hệ kín có tổng ngoại lực bằng 0 hoặc không đáng kể. (Vì): Phù hợp với định nghĩa hoặc hệ thức vật lí. b) Đúng: Phù hợp với định nghĩa hoặc hệ thức vật lí. c) Sai: Không phù hợp với định nghĩa hoặc hệ thức vật lí. d) Đúng: Phù hợp với định nghĩa hoặc hệ thức vật lí
\item (Đúng) Trong hệ kín, tổng động lượng được bảo toàn.
\item (Sai) Mọi hệ vật đều là hệ kín.
\item (Đúng) Nội lực không làm thay đổi tổng động lượng hệ.
\end{itemchoice}
}
\end{ex}

% ===== Câu 12 =====
% ID: AIQ_L10C5_FULL_0174
% Mức: TH
\begin{ex}
% ID: AIQ_L10C5_FULL_0174
% Mức: TH
Xét các phát biểu thuộc dạng “Nhận biết hệ kín và điều kiện bảo toàn động lượng” ở tình huống số 2.
\choiceTF
{\True Hệ kín có tổng ngoại lực bằng 0 hoặc không đáng kể.}
{\True Trong hệ kín, tổng động lượng được bảo toàn.}
{Mọi hệ vật đều là hệ kín.}
{\True Nội lực không làm thay đổi tổng động lượng hệ.}
\loigiai{
\begin{itemchoice}
\item (Đúng) Hệ kín có tổng ngoại lực bằng 0 hoặc không đáng kể. (Vì): Phù hợp với định nghĩa hoặc hệ thức vật lí. b) Đúng: Phù hợp với định nghĩa hoặc hệ thức vật lí. c) Sai: Không phù hợp với định nghĩa hoặc hệ thức vật lí. d) Đúng: Phù hợp với định nghĩa hoặc hệ thức vật lí
\item (Đúng) Trong hệ kín, tổng động lượng được bảo toàn.
\item (Sai) Mọi hệ vật đều là hệ kín.
\item (Đúng) Nội lực không làm thay đổi tổng động lượng hệ.
\end{itemchoice}
}
\end{ex}

% ===== Câu 13 =====
% ID: AIQ_L10C5_FULL_0175
% Mức: TH
\begin{ex}
% ID: AIQ_L10C5_FULL_0175
% Mức: TH
Xét các phát biểu thuộc dạng “Nhận biết hệ kín và điều kiện bảo toàn động lượng” ở tình huống số 3.
\choiceTF
{\True Hệ kín có tổng ngoại lực bằng 0 hoặc không đáng kể.}
{\True Trong hệ kín, tổng động lượng được bảo toàn.}
{Mọi hệ vật đều là hệ kín.}
{\True Nội lực không làm thay đổi tổng động lượng hệ.}
\loigiai{
\begin{itemchoice}
\item (Đúng) Hệ kín có tổng ngoại lực bằng 0 hoặc không đáng kể. (Vì): Phù hợp với định nghĩa hoặc hệ thức vật lí. b) Đúng: Phù hợp với định nghĩa hoặc hệ thức vật lí. c) Sai: Không phù hợp với định nghĩa hoặc hệ thức vật lí. d) Đúng: Phù hợp với định nghĩa hoặc hệ thức vật lí
\item (Đúng) Trong hệ kín, tổng động lượng được bảo toàn.
\item (Sai) Mọi hệ vật đều là hệ kín.
\item (Đúng) Nội lực không làm thay đổi tổng động lượng hệ.
\end{itemchoice}
}
\end{ex}

% ===== Câu 14 =====
% ID: AIQ_L10C5_FULL_0176
% Mức: VD
\begin{ex}
% ID: AIQ_L10C5_FULL_0176
% Mức: VD
Xét các phát biểu thuộc dạng “Nhận biết hệ kín và điều kiện bảo toàn động lượng” ở tình huống số 4.
\choiceTF
{\True Hệ kín có tổng ngoại lực bằng 0 hoặc không đáng kể.}
{\True Trong hệ kín, tổng động lượng được bảo toàn.}
{Mọi hệ vật đều là hệ kín.}
{\True Nội lực không làm thay đổi tổng động lượng hệ.}
\loigiai{
\begin{itemchoice}
\item (Đúng) Hệ kín có tổng ngoại lực bằng 0 hoặc không đáng kể. (Vì): Phù hợp với định nghĩa hoặc hệ thức vật lí. b) Đúng: Phù hợp với định nghĩa hoặc hệ thức vật lí. c) Sai: Không phù hợp với định nghĩa hoặc hệ thức vật lí. d) Đúng: Phù hợp với định nghĩa hoặc hệ thức vật lí
\item (Đúng) Trong hệ kín, tổng động lượng được bảo toàn.
\item (Sai) Mọi hệ vật đều là hệ kín.
\item (Đúng) Nội lực không làm thay đổi tổng động lượng hệ.
\end{itemchoice}
}
\end{ex}

% ===== Câu 15 =====
% ID: AIQ_L10C5_FULL_0177
% Mức: VDC
\begin{ex}
% ID: AIQ_L10C5_FULL_0177
% Mức: VDC
Xét các phát biểu thuộc dạng “Nhận biết hệ kín và điều kiện bảo toàn động lượng” ở tình huống số 5.
\choiceTF
{\True Hệ kín có tổng ngoại lực bằng 0 hoặc không đáng kể.}
{\True Trong hệ kín, tổng động lượng được bảo toàn.}
{Mọi hệ vật đều là hệ kín.}
{\True Nội lực không làm thay đổi tổng động lượng hệ.}
\loigiai{
\begin{itemchoice}
\item (Đúng) Hệ kín có tổng ngoại lực bằng 0 hoặc không đáng kể. (Vì): Phù hợp với định nghĩa hoặc hệ thức vật lí. b) Đúng: Phù hợp với định nghĩa hoặc hệ thức vật lí. c) Sai: Không phù hợp với định nghĩa hoặc hệ thức vật lí. d) Đúng: Phù hợp với định nghĩa hoặc hệ thức vật lí
\item (Đúng) Trong hệ kín, tổng động lượng được bảo toàn.
\item (Sai) Mọi hệ vật đều là hệ kín.
\item (Đúng) Nội lực không làm thay đổi tổng động lượng hệ.
\end{itemchoice}
}
\end{ex}

% ===== Câu 16 =====
% ID: AIQ_L10C5_FULL_0178
% Mức: TH
\begin{ex}
% ID: AIQ_L10C5_FULL_0178
% Mức: TH
Một hệ kín ban đầu đứng yên tách thành hai phần. Phần thứ nhất có động lượng $52\ \mathrm{kg\,m/s}$ theo chiều dương. Độ lớn động lượng phần thứ hai bằng bao nhiêu?
\shortans{52}
\loigiai{
Hệ kín có tổng động lượng ban đầu bằng 0 nên $\vec p_2=-\vec p_1$. Suy ra kết quả $52\ \mathrm{kg\,m/s}$. Đáp án trả lời ngắn không ghi đơn vị.
}
\end{ex}

% ===== Câu 17 =====
% ID: AIQ_L10C5_FULL_0179
% Mức: TH
\begin{ex}
% ID: AIQ_L10C5_FULL_0179
% Mức: TH
Một hệ kín ban đầu đứng yên tách thành hai phần. Phần thứ nhất có động lượng $70\ \mathrm{kg\,m/s}$ theo chiều dương. Độ lớn động lượng phần thứ hai bằng bao nhiêu?
\shortans{70}
\loigiai{
Hệ kín có tổng động lượng ban đầu bằng 0 nên $\vec p_2=-\vec p_1$. Suy ra kết quả $70\ \mathrm{kg\,m/s}$. Đáp án trả lời ngắn không ghi đơn vị.
}
\end{ex}

% ===== Câu 18 =====
% ID: AIQ_L10C5_FULL_0180
% Mức: VD
\begin{ex}
% ID: AIQ_L10C5_FULL_0180
% Mức: VD
Một hệ kín ban đầu đứng yên tách thành hai phần. Phần thứ nhất có động lượng $30\ \mathrm{kg\,m/s}$ theo chiều dương. Độ lớn động lượng phần thứ hai bằng bao nhiêu?
\shortans{30}
\loigiai{
Hệ kín có tổng động lượng ban đầu bằng 0 nên $\vec p_2=-\vec p_1$. Suy ra kết quả $30\ \mathrm{kg\,m/s}$. Đáp án trả lời ngắn không ghi đơn vị.
}
\end{ex}

% ===== Câu 19 =====
% ID: AIQ_L10C5_FULL_0181
% Mức: VD
\begin{ex}
% ID: AIQ_L10C5_FULL_0181
% Mức: VD
Một hệ kín ban đầu đứng yên tách thành hai phần. Phần thứ nhất có động lượng $48\ \mathrm{kg\,m/s}$ theo chiều dương. Độ lớn động lượng phần thứ hai bằng bao nhiêu?
\shortans{48}
\loigiai{
Hệ kín có tổng động lượng ban đầu bằng 0 nên $\vec p_2=-\vec p_1$. Suy ra kết quả $48\ \mathrm{kg\,m/s}$. Đáp án trả lời ngắn không ghi đơn vị.
}
\end{ex}

% ===== Câu 20 =====
% ID: AIQ_L10C5_FULL_0182
% Mức: VDC
\begin{ex}
% ID: AIQ_L10C5_FULL_0182
% Mức: VDC
Một hệ kín ban đầu đứng yên tách thành hai phần. Phần thứ nhất có động lượng $68\ \mathrm{kg\,m/s}$ theo chiều dương. Độ lớn động lượng phần thứ hai bằng bao nhiêu?
\shortans{68}
\loigiai{
Hệ kín có tổng động lượng ban đầu bằng 0 nên $\vec p_2=-\vec p_1$. Suy ra kết quả $68\ \mathrm{kg\,m/s}$. Đáp án trả lời ngắn không ghi đơn vị.
}
\end{ex}

% ===== Câu 21 =====
% ID: AIQ_L10C5_FULL_0183
% Mức: VDC
\begin{ex}
% ID: AIQ_L10C5_FULL_0183
% Mức: VDC
Một hệ kín ban đầu đứng yên tách thành hai phần. Phần thứ nhất có động lượng $90\ \mathrm{kg\,m/s}$ theo chiều dương. Độ lớn động lượng phần thứ hai bằng bao nhiêu?
\shortans{90}
\loigiai{
Hệ kín có tổng động lượng ban đầu bằng 0 nên $\vec p_2=-\vec p_1$. Suy ra kết quả $90\ \mathrm{kg\,m/s}$. Đáp án trả lời ngắn không ghi đơn vị.
}
\end{ex}

% ===== Câu 22 =====
% ID: AIQ_L10C5_FULL_0184
% Mức: TH
\begin{bt}
% ID: AIQ_L10C5_FULL_0184
% Mức: TH
Trình bày cách giải: Một hệ kín ban đầu đứng yên tách thành hai phần. Phần thứ nhất có động lượng $38\ \mathrm{kg\,m/s}$ theo chiều dương. Độ lớn động lượng phần thứ hai bằng bao nhiêu?
\loigiai{
Hệ kín có tổng động lượng ban đầu bằng 0 nên $\vec p_2=-\vec p_1$. Suy ra kết quả $38\ \mathrm{kg\,m/s}$.
}
\end{bt}

% ===== Câu 23 =====
% ID: AIQ_L10C5_FULL_0185
% Mức: TH
\begin{bt}
% ID: AIQ_L10C5_FULL_0185
% Mức: TH
Trình bày cách giải: Một hệ kín ban đầu đứng yên tách thành hai phần. Phần thứ nhất có động lượng $60\ \mathrm{kg\,m/s}$ theo chiều dương. Độ lớn động lượng phần thứ hai bằng bao nhiêu?
\loigiai{
Hệ kín có tổng động lượng ban đầu bằng 0 nên $\vec p_2=-\vec p_1$. Suy ra kết quả $60\ \mathrm{kg\,m/s}$.
}
\end{bt}

% ===== Câu 24 =====
% ID: AIQ_L10C5_FULL_0186
% Mức: VD
\begin{bt}
% ID: AIQ_L10C5_FULL_0186
% Mức: VD
Trình bày cách giải: Một hệ kín ban đầu đứng yên tách thành hai phần. Phần thứ nhất có động lượng $84\ \mathrm{kg\,m/s}$ theo chiều dương. Độ lớn động lượng phần thứ hai bằng bao nhiêu?
\loigiai{
Hệ kín có tổng động lượng ban đầu bằng 0 nên $\vec p_2=-\vec p_1$. Suy ra kết quả $84\ \mathrm{kg\,m/s}$.
}
\end{bt}

% ===== Câu 25 =====
% ID: AIQ_L10C5_FULL_0187
% Mức: VD
\begin{bt}
% ID: AIQ_L10C5_FULL_0187
% Mức: VD
Trình bày cách giải: Một hệ kín ban đầu đứng yên tách thành hai phần. Phần thứ nhất có động lượng $110\ \mathrm{kg\,m/s}$ theo chiều dương. Độ lớn động lượng phần thứ hai bằng bao nhiêu?
\loigiai{
Hệ kín có tổng động lượng ban đầu bằng 0 nên $\vec p_2=-\vec p_1$. Suy ra kết quả $110\ \mathrm{kg\,m/s}$.
}
\end{bt}

% ===== Câu 26 =====
% ID: AIQ_L10C5_FULL_0188
% Mức: VDC
\begin{bt}
% ID: AIQ_L10C5_FULL_0188
% Mức: VDC
Trình bày cách giải: Một hệ kín ban đầu đứng yên tách thành hai phần. Phần thứ nhất có động lượng $46\ \mathrm{kg\,m/s}$ theo chiều dương. Độ lớn động lượng phần thứ hai bằng bao nhiêu?
\loigiai{
Hệ kín có tổng động lượng ban đầu bằng 0 nên $\vec p_2=-\vec p_1$. Suy ra kết quả $46\ \mathrm{kg\,m/s}$.
}
\end{bt}

% ===== Câu 27 =====
% ID: AIQ_L10C5_FULL_0189
% Mức: VDC
\begin{bt}
% ID: AIQ_L10C5_FULL_0189
% Mức: VDC
Trình bày cách giải: Một hệ kín ban đầu đứng yên tách thành hai phần. Phần thứ nhất có động lượng $72\ \mathrm{kg\,m/s}$ theo chiều dương. Độ lớn động lượng phần thứ hai bằng bao nhiêu?
\loigiai{
Hệ kín có tổng động lượng ban đầu bằng 0 nên $\vec p_2=-\vec p_1$. Suy ra kết quả $72\ \mathrm{kg\,m/s}$.
}
\end{bt}

% ===== Câu 28 =====
% ID: AIQ_L10C5_FULL_0190
% Mức: NB
\dangbt{Bảo toàn động lượng trong va chạm thẳng}
\begin{ex}
% ID: AIQ_L10C5_FULL_0190
% Mức: NB
Chọn chiều dương là chiều vật 1. Trước va chạm: $m_1=1$ kg, $v_1=4$ m/s; $m_2=2$ kg, $v_2=-1$ m/s. Sau va chạm vật 2 đứng yên. Tính vận tốc vật 1.
\choice
{\True $2\ \mathrm{m/s}$}
{Không đủ dữ kiện để có giá trị này}
{$4\ \mathrm{m/s}$}
{$4\ \mathrm{m/s}$}
\loigiai{
Bảo toàn động lượng: $m_1v_1+m_2v_2=m_1v'_1$, nên $v'_1=(1\cdot4-2\cdot1)/1$. Suy ra kết quả $2\ \mathrm{m/s}$.
}
\end{ex}

% ===== Câu 29 =====
% ID: AIQ_L10C5_FULL_0191
% Mức: NB
\begin{ex}
% ID: AIQ_L10C5_FULL_0191
% Mức: NB
Chọn chiều dương là chiều vật 1. Trước va chạm: $m_1=2$ kg, $v_1=5$ m/s; $m_2=3$ kg, $v_2=-2$ m/s. Sau va chạm vật 2 đứng yên. Tính vận tốc vật 1.
\choice
{Không đủ dữ kiện để có giá trị này}
{\True $2\ \mathrm{m/s}$}
{$4\ \mathrm{m/s}$}
{$4\ \mathrm{m/s}$}
\loigiai{
Bảo toàn động lượng: $m_1v_1+m_2v_2=m_1v'_1$, nên $v'_1=(2\cdot5-3\cdot2)/2$. Suy ra kết quả $2\ \mathrm{m/s}$.
}
\end{ex}

% ===== Câu 30 =====
% ID: AIQ_L10C5_FULL_0192
% Mức: TH
\begin{ex}
% ID: AIQ_L10C5_FULL_0192
% Mức: TH
Chọn chiều dương là chiều vật 1. Trước va chạm: $m_1=3$ kg, $v_1=6$ m/s; $m_2=4$ kg, $v_2=-1$ m/s. Sau va chạm vật 2 đứng yên. Tính vận tốc vật 1.
\choice
{Không đủ dữ kiện để có giá trị này}
{$6,67\ \mathrm{m/s}$}
{\True $4,67\ \mathrm{m/s}$}
{$9,33\ \mathrm{m/s}$}
\loigiai{
Bảo toàn động lượng: $m_1v_1+m_2v_2=m_1v'_1$, nên $v'_1=(3\cdot6-4\cdot1)/3$. Suy ra kết quả $4,67\ \mathrm{m/s}$.
}
\end{ex}

% ===== Câu 31 =====
% ID: AIQ_L10C5_FULL_0193
% Mức: TH
\begin{ex}
% ID: AIQ_L10C5_FULL_0193
% Mức: TH
Chọn chiều dương là chiều vật 1. Trước va chạm: $m_1=1$ kg, $v_1=7$ m/s; $m_2=5$ kg, $v_2=-2$ m/s. Sau va chạm vật 2 đứng yên. Tính vận tốc vật 1.
\choice
{$3\ \mathrm{m/s}$}
{$-1\ \mathrm{m/s}$}
{$-6\ \mathrm{m/s}$}
{\True $-3\ \mathrm{m/s}$}
\loigiai{
Bảo toàn động lượng: $m_1v_1+m_2v_2=m_1v'_1$, nên $v'_1=(1\cdot7-5\cdot2)/1$. Suy ra kết quả $-3\ \mathrm{m/s}$.
}
\end{ex}

% ===== Câu 32 =====
% ID: AIQ_L10C5_FULL_0194
% Mức: VD
\begin{ex}
% ID: AIQ_L10C5_FULL_0194
% Mức: VD
Chọn chiều dương là chiều vật 1. Trước va chạm: $m_1=2$ kg, $v_1=8$ m/s; $m_2=2$ kg, $v_2=-1$ m/s. Sau va chạm vật 2 đứng yên. Tính vận tốc vật 1.
\choice
{\True $7\ \mathrm{m/s}$}
{Không đủ dữ kiện để có giá trị này}
{$9\ \mathrm{m/s}$}
{$14\ \mathrm{m/s}$}
\loigiai{
Bảo toàn động lượng: $m_1v_1+m_2v_2=m_1v'_1$, nên $v'_1=(2\cdot8-2\cdot1)/2$. Suy ra kết quả $7\ \mathrm{m/s}$.
}
\end{ex}

% ===== Câu 33 =====
% ID: AIQ_L10C5_FULL_0195
% Mức: VD
\begin{ex}
% ID: AIQ_L10C5_FULL_0195
% Mức: VD
Chọn chiều dương là chiều vật 1. Trước va chạm: $m_1=3$ kg, $v_1=9$ m/s; $m_2=3$ kg, $v_2=-2$ m/s. Sau va chạm vật 2 đứng yên. Tính vận tốc vật 1.
\choice
{Không đủ dữ kiện để có giá trị này}
{\True $7\ \mathrm{m/s}$}
{$9\ \mathrm{m/s}$}
{$14\ \mathrm{m/s}$}
\loigiai{
Bảo toàn động lượng: $m_1v_1+m_2v_2=m_1v'_1$, nên $v'_1=(3\cdot9-3\cdot2)/3$. Suy ra kết quả $7\ \mathrm{m/s}$.
}
\end{ex}

% ===== Câu 34 =====
% ID: AIQ_L10C5_FULL_0196
% Mức: VD
\begin{ex}
% ID: AIQ_L10C5_FULL_0196
% Mức: VD
Chọn chiều dương là chiều vật 1. Trước va chạm: $m_1=1$ kg, $v_1=10$ m/s; $m_2=4$ kg, $v_2=-1$ m/s. Sau va chạm vật 2 đứng yên. Tính vận tốc vật 1.
\choice
{Không đủ dữ kiện để có giá trị này}
{$8\ \mathrm{m/s}$}
{\True $6\ \mathrm{m/s}$}
{$12\ \mathrm{m/s}$}
\loigiai{
Bảo toàn động lượng: $m_1v_1+m_2v_2=m_1v'_1$, nên $v'_1=(1\cdot10-4\cdot1)/1$. Suy ra kết quả $6\ \mathrm{m/s}$.
}
\end{ex}

% ===== Câu 35 =====
% ID: AIQ_L10C5_FULL_0197
% Mức: VDC
\begin{ex}
% ID: AIQ_L10C5_FULL_0197
% Mức: VDC
Chọn chiều dương là chiều vật 1. Trước va chạm: $m_1=2$ kg, $v_1=11$ m/s; $m_2=5$ kg, $v_2=-2$ m/s. Sau va chạm vật 2 đứng yên. Tính vận tốc vật 1.
\choice
{Không đủ dữ kiện để có giá trị này}
{$8\ \mathrm{m/s}$}
{$12\ \mathrm{m/s}$}
{\True $6\ \mathrm{m/s}$}
\loigiai{
Bảo toàn động lượng: $m_1v_1+m_2v_2=m_1v'_1$, nên $v'_1=(2\cdot11-5\cdot2)/2$. Suy ra kết quả $6\ \mathrm{m/s}$.
}
\end{ex}

% ===== Câu 36 =====
% ID: AIQ_L10C5_FULL_0198
% Mức: VDC
\begin{ex}
% ID: AIQ_L10C5_FULL_0198
% Mức: VDC
Chọn chiều dương là chiều vật 1. Trước va chạm: $m_1=3$ kg, $v_1=12$ m/s; $m_2=2$ kg, $v_2=-1$ m/s. Sau va chạm vật 2 đứng yên. Tính vận tốc vật 1.
\choice
{\True $11,33\ \mathrm{m/s}$}
{Không đủ dữ kiện để có giá trị này}
{$13,33\ \mathrm{m/s}$}
{$22,67\ \mathrm{m/s}$}
\loigiai{
Bảo toàn động lượng: $m_1v_1+m_2v_2=m_1v'_1$, nên $v'_1=(3\cdot12-2\cdot1)/3$. Suy ra kết quả $11,33\ \mathrm{m/s}$.
}
\end{ex}

% ===== Câu 37 =====
% ID: AIQ_L10C5_FULL_0199
% Mức: VDC
\begin{ex}
% ID: AIQ_L10C5_FULL_0199
% Mức: VDC
Chọn chiều dương là chiều vật 1. Trước va chạm: $m_1=1$ kg, $v_1=13$ m/s; $m_2=3$ kg, $v_2=-2$ m/s. Sau va chạm vật 2 đứng yên. Tính vận tốc vật 1.
\choice
{Không đủ dữ kiện để có giá trị này}
{\True $7\ \mathrm{m/s}$}
{$9\ \mathrm{m/s}$}
{$14\ \mathrm{m/s}$}
\loigiai{
Bảo toàn động lượng: $m_1v_1+m_2v_2=m_1v'_1$, nên $v'_1=(1\cdot13-3\cdot2)/1$. Suy ra kết quả $7\ \mathrm{m/s}$.
}
\end{ex}

% ===== Câu 38 =====
% ID: AIQ_L10C5_FULL_0200
% Mức: NB
\begin{ex}
% ID: AIQ_L10C5_FULL_0200
% Mức: NB
Xét các phát biểu thuộc dạng “Bảo toàn động lượng trong va chạm thẳng” ở tình huống số 1.
\choiceTF
{\True Va chạm thẳng được xét trên một trục.}
{\True Phải chọn chiều dương khi dùng đại số.}
{Động lượng từng vật luôn bảo toàn.}
{\True Tổng động lượng hệ kín bảo toàn.}
\loigiai{
\begin{itemchoice}
\item (Đúng) Va chạm thẳng được xét trên một trục. (Vì): Phù hợp với định nghĩa hoặc hệ thức vật lí. b) Đúng: Phù hợp với định nghĩa hoặc hệ thức vật lí. c) Sai: Không phù hợp với định nghĩa hoặc hệ thức vật lí. d) Đúng: Phù hợp với định nghĩa hoặc hệ thức vật lí
\item (Đúng) Phải chọn chiều dương khi dùng đại số.
\item (Sai) Động lượng từng vật luôn bảo toàn.
\item (Đúng) Tổng động lượng hệ kín bảo toàn.
\end{itemchoice}
}
\end{ex}

% ===== Câu 39 =====
% ID: AIQ_L10C5_FULL_0201
% Mức: TH
\begin{ex}
% ID: AIQ_L10C5_FULL_0201
% Mức: TH
Xét các phát biểu thuộc dạng “Bảo toàn động lượng trong va chạm thẳng” ở tình huống số 2.
\choiceTF
{\True Va chạm thẳng được xét trên một trục.}
{\True Phải chọn chiều dương khi dùng đại số.}
{Động lượng từng vật luôn bảo toàn.}
{\True Tổng động lượng hệ kín bảo toàn.}
\loigiai{
\begin{itemchoice}
\item (Đúng) Va chạm thẳng được xét trên một trục. (Vì): Phù hợp với định nghĩa hoặc hệ thức vật lí. b) Đúng: Phù hợp với định nghĩa hoặc hệ thức vật lí. c) Sai: Không phù hợp với định nghĩa hoặc hệ thức vật lí. d) Đúng: Phù hợp với định nghĩa hoặc hệ thức vật lí
\item (Đúng) Phải chọn chiều dương khi dùng đại số.
\item (Sai) Động lượng từng vật luôn bảo toàn.
\item (Đúng) Tổng động lượng hệ kín bảo toàn.
\end{itemchoice}
}
\end{ex}

% ===== Câu 40 =====
% ID: AIQ_L10C5_FULL_0202
% Mức: TH
\begin{ex}
% ID: AIQ_L10C5_FULL_0202
% Mức: TH
Xét các phát biểu thuộc dạng “Bảo toàn động lượng trong va chạm thẳng” ở tình huống số 3.
\choiceTF
{\True Va chạm thẳng được xét trên một trục.}
{\True Phải chọn chiều dương khi dùng đại số.}
{Động lượng từng vật luôn bảo toàn.}
{\True Tổng động lượng hệ kín bảo toàn.}
\loigiai{
\begin{itemchoice}
\item (Đúng) Va chạm thẳng được xét trên một trục. (Vì): Phù hợp với định nghĩa hoặc hệ thức vật lí. b) Đúng: Phù hợp với định nghĩa hoặc hệ thức vật lí. c) Sai: Không phù hợp với định nghĩa hoặc hệ thức vật lí. d) Đúng: Phù hợp với định nghĩa hoặc hệ thức vật lí
\item (Đúng) Phải chọn chiều dương khi dùng đại số.
\item (Sai) Động lượng từng vật luôn bảo toàn.
\item (Đúng) Tổng động lượng hệ kín bảo toàn.
\end{itemchoice}
}
\end{ex}

% ===== Câu 41 =====
% ID: AIQ_L10C5_FULL_0203
% Mức: VD
\begin{ex}
% ID: AIQ_L10C5_FULL_0203
% Mức: VD
Xét các phát biểu thuộc dạng “Bảo toàn động lượng trong va chạm thẳng” ở tình huống số 4.
\choiceTF
{\True Va chạm thẳng được xét trên một trục.}
{\True Phải chọn chiều dương khi dùng đại số.}
{Động lượng từng vật luôn bảo toàn.}
{\True Tổng động lượng hệ kín bảo toàn.}
\loigiai{
\begin{itemchoice}
\item (Đúng) Va chạm thẳng được xét trên một trục. (Vì): Phù hợp với định nghĩa hoặc hệ thức vật lí. b) Đúng: Phù hợp với định nghĩa hoặc hệ thức vật lí. c) Sai: Không phù hợp với định nghĩa hoặc hệ thức vật lí. d) Đúng: Phù hợp với định nghĩa hoặc hệ thức vật lí
\item (Đúng) Phải chọn chiều dương khi dùng đại số.
\item (Sai) Động lượng từng vật luôn bảo toàn.
\item (Đúng) Tổng động lượng hệ kín bảo toàn.
\end{itemchoice}
}
\end{ex}

% ===== Câu 42 =====
% ID: AIQ_L10C5_FULL_0204
% Mức: VDC
\begin{ex}
% ID: AIQ_L10C5_FULL_0204
% Mức: VDC
Xét các phát biểu thuộc dạng “Bảo toàn động lượng trong va chạm thẳng” ở tình huống số 5.
\choiceTF
{\True Va chạm thẳng được xét trên một trục.}
{\True Phải chọn chiều dương khi dùng đại số.}
{Động lượng từng vật luôn bảo toàn.}
{\True Tổng động lượng hệ kín bảo toàn.}
\loigiai{
\begin{itemchoice}
\item (Đúng) Va chạm thẳng được xét trên một trục. (Vì): Phù hợp với định nghĩa hoặc hệ thức vật lí. b) Đúng: Phù hợp với định nghĩa hoặc hệ thức vật lí. c) Sai: Không phù hợp với định nghĩa hoặc hệ thức vật lí. d) Đúng: Phù hợp với định nghĩa hoặc hệ thức vật lí
\item (Đúng) Phải chọn chiều dương khi dùng đại số.
\item (Sai) Động lượng từng vật luôn bảo toàn.
\item (Đúng) Tổng động lượng hệ kín bảo toàn.
\end{itemchoice}
}
\end{ex}

% ===== Câu 43 =====
% ID: AIQ_L10C5_FULL_0205
% Mức: TH
\begin{ex}
% ID: AIQ_L10C5_FULL_0205
% Mức: TH
Chọn chiều dương là chiều vật 1. Trước va chạm: $m_1=2$ kg, $v_1=14$ m/s; $m_2=4$ kg, $v_2=-1$ m/s. Sau va chạm vật 2 đứng yên. Tính vận tốc vật 1.
\shortans{12}
\loigiai{
Bảo toàn động lượng: $m_1v_1+m_2v_2=m_1v'_1$, nên $v'_1=(2\cdot14-4\cdot1)/2$. Suy ra kết quả $12\ \mathrm{m/s}$. Đáp án trả lời ngắn không ghi đơn vị.
}
\end{ex}

% ===== Câu 44 =====
% ID: AIQ_L10C5_FULL_0206
% Mức: TH
\begin{ex}
% ID: AIQ_L10C5_FULL_0206
% Mức: TH
Chọn chiều dương là chiều vật 1. Trước va chạm: $m_1=3$ kg, $v_1=15$ m/s; $m_2=5$ kg, $v_2=-2$ m/s. Sau va chạm vật 2 đứng yên. Tính vận tốc vật 1.
\shortans{11,67}
\loigiai{
Bảo toàn động lượng: $m_1v_1+m_2v_2=m_1v'_1$, nên $v'_1=(3\cdot15-5\cdot2)/3$. Suy ra kết quả $11,67\ \mathrm{m/s}$. Đáp án trả lời ngắn không ghi đơn vị.
}
\end{ex}

% ===== Câu 45 =====
% ID: AIQ_L10C5_FULL_0207
% Mức: VD
\begin{ex}
% ID: AIQ_L10C5_FULL_0207
% Mức: VD
Chọn chiều dương là chiều vật 1. Trước va chạm: $m_1=1$ kg, $v_1=16$ m/s; $m_2=2$ kg, $v_2=-1$ m/s. Sau va chạm vật 2 đứng yên. Tính vận tốc vật 1.
\shortans{14}
\loigiai{
Bảo toàn động lượng: $m_1v_1+m_2v_2=m_1v'_1$, nên $v'_1=(1\cdot16-2\cdot1)/1$. Suy ra kết quả $14\ \mathrm{m/s}$. Đáp án trả lời ngắn không ghi đơn vị.
}
\end{ex}

% ===== Câu 46 =====
% ID: AIQ_L10C5_FULL_0208
% Mức: VD
\begin{ex}
% ID: AIQ_L10C5_FULL_0208
% Mức: VD
Chọn chiều dương là chiều vật 1. Trước va chạm: $m_1=2$ kg, $v_1=17$ m/s; $m_2=3$ kg, $v_2=-2$ m/s. Sau va chạm vật 2 đứng yên. Tính vận tốc vật 1.
\shortans{14}
\loigiai{
Bảo toàn động lượng: $m_1v_1+m_2v_2=m_1v'_1$, nên $v'_1=(2\cdot17-3\cdot2)/2$. Suy ra kết quả $14\ \mathrm{m/s}$. Đáp án trả lời ngắn không ghi đơn vị.
}
\end{ex}

% ===== Câu 47 =====
% ID: AIQ_L10C5_FULL_0209
% Mức: VDC
\begin{ex}
% ID: AIQ_L10C5_FULL_0209
% Mức: VDC
Chọn chiều dương là chiều vật 1. Trước va chạm: $m_1=3$ kg, $v_1=18$ m/s; $m_2=4$ kg, $v_2=-1$ m/s. Sau va chạm vật 2 đứng yên. Tính vận tốc vật 1.
\shortans{16,67}
\loigiai{
Bảo toàn động lượng: $m_1v_1+m_2v_2=m_1v'_1$, nên $v'_1=(3\cdot18-4\cdot1)/3$. Suy ra kết quả $16,67\ \mathrm{m/s}$. Đáp án trả lời ngắn không ghi đơn vị.
}
\end{ex}

% ===== Câu 48 =====
% ID: AIQ_L10C5_FULL_0210
% Mức: VDC
\begin{ex}
% ID: AIQ_L10C5_FULL_0210
% Mức: VDC
Chọn chiều dương là chiều vật 1. Trước va chạm: $m_1=1$ kg, $v_1=19$ m/s; $m_2=5$ kg, $v_2=-2$ m/s. Sau va chạm vật 2 đứng yên. Tính vận tốc vật 1.
\shortans{9}
\loigiai{
Bảo toàn động lượng: $m_1v_1+m_2v_2=m_1v'_1$, nên $v'_1=(1\cdot19-5\cdot2)/1$. Suy ra kết quả $9\ \mathrm{m/s}$. Đáp án trả lời ngắn không ghi đơn vị.
}
\end{ex}

% ===== Câu 49 =====
% ID: AIQ_L10C5_FULL_0211
% Mức: TH
\begin{bt}
% ID: AIQ_L10C5_FULL_0211
% Mức: TH
Trình bày cách giải: Chọn chiều dương là chiều vật 1. Trước va chạm: $m_1=2$ kg, $v_1=20$ m/s; $m_2=2$ kg, $v_2=-1$ m/s. Sau va chạm vật 2 đứng yên. Tính vận tốc vật 1.
\loigiai{
Bảo toàn động lượng: $m_1v_1+m_2v_2=m_1v'_1$, nên $v'_1=(2\cdot20-2\cdot1)/2$. Suy ra kết quả $19\ \mathrm{m/s}$.
}
\end{bt}

% ===== Câu 50 =====
% ID: AIQ_L10C5_FULL_0212
% Mức: TH
\begin{bt}
% ID: AIQ_L10C5_FULL_0212
% Mức: TH
Trình bày cách giải: Chọn chiều dương là chiều vật 1. Trước va chạm: $m_1=3$ kg, $v_1=21$ m/s; $m_2=3$ kg, $v_2=-2$ m/s. Sau va chạm vật 2 đứng yên. Tính vận tốc vật 1.
\loigiai{
Bảo toàn động lượng: $m_1v_1+m_2v_2=m_1v'_1$, nên $v'_1=(3\cdot21-3\cdot2)/3$. Suy ra kết quả $19\ \mathrm{m/s}$.
}
\end{bt}

% ===== Câu 51 =====
% ID: AIQ_L10C5_FULL_0213
% Mức: VD
\begin{bt}
% ID: AIQ_L10C5_FULL_0213
% Mức: VD
Trình bày cách giải: Chọn chiều dương là chiều vật 1. Trước va chạm: $m_1=1$ kg, $v_1=22$ m/s; $m_2=4$ kg, $v_2=-1$ m/s. Sau va chạm vật 2 đứng yên. Tính vận tốc vật 1.
\loigiai{
Bảo toàn động lượng: $m_1v_1+m_2v_2=m_1v'_1$, nên $v'_1=(1\cdot22-4\cdot1)/1$. Suy ra kết quả $18\ \mathrm{m/s}$.
}
\end{bt}

% ===== Câu 52 =====
% ID: AIQ_L10C5_FULL_0214
% Mức: VD
\begin{bt}
% ID: AIQ_L10C5_FULL_0214
% Mức: VD
Trình bày cách giải: Chọn chiều dương là chiều vật 1. Trước va chạm: $m_1=2$ kg, $v_1=23$ m/s; $m_2=5$ kg, $v_2=-2$ m/s. Sau va chạm vật 2 đứng yên. Tính vận tốc vật 1.
\loigiai{
Bảo toàn động lượng: $m_1v_1+m_2v_2=m_1v'_1$, nên $v'_1=(2\cdot23-5\cdot2)/2$. Suy ra kết quả $18\ \mathrm{m/s}$.
}
\end{bt}

% ===== Câu 53 =====
% ID: AIQ_L10C5_FULL_0215
% Mức: VDC
\begin{bt}
% ID: AIQ_L10C5_FULL_0215
% Mức: VDC
Trình bày cách giải: Chọn chiều dương là chiều vật 1. Trước va chạm: $m_1=3$ kg, $v_1=24$ m/s; $m_2=2$ kg, $v_2=-1$ m/s. Sau va chạm vật 2 đứng yên. Tính vận tốc vật 1.
\loigiai{
Bảo toàn động lượng: $m_1v_1+m_2v_2=m_1v'_1$, nên $v'_1=(3\cdot24-2\cdot1)/3$. Suy ra kết quả $23,33\ \mathrm{m/s}$.
}
\end{bt}

% ===== Câu 54 =====
% ID: AIQ_L10C5_FULL_0216
% Mức: VDC
\begin{bt}
% ID: AIQ_L10C5_FULL_0216
% Mức: VDC
Trình bày cách giải: Chọn chiều dương là chiều vật 1. Trước va chạm: $m_1=1$ kg, $v_1=25$ m/s; $m_2=3$ kg, $v_2=-2$ m/s. Sau va chạm vật 2 đứng yên. Tính vận tốc vật 1.
\loigiai{
Bảo toàn động lượng: $m_1v_1+m_2v_2=m_1v'_1$, nên $v'_1=(1\cdot25-3\cdot2)/1$. Suy ra kết quả $19\ \mathrm{m/s}$.
}
\end{bt}

% ===== Câu 55 =====
% ID: AIQ_L10C5_FULL_0217
% Mức: NB
\dangbt{Bảo toàn động lượng trong va chạm mềm}
\begin{ex}
% ID: AIQ_L10C5_FULL_0217
% Mức: NB
Vật $m_1=1$ kg chuyển động $4$ m/s va chạm mềm cùng chiều với vật $m_2=2$ kg chuyển động $0$ m/s. Tính vận tốc chung.
\choice
{\True $1,33\ \mathrm{m/s}$}
{$4\ \mathrm{m/s}$}
{$4\ \mathrm{m/s}$}
{$4\ \mathrm{m/s}$}
\loigiai{
$v'=\dfrac{m_1v_1+m_2v_2}{m_1+m_2}=\dfrac{1\cdot4+2\cdot0}{1+2}$. Suy ra kết quả $1,33\ \mathrm{m/s}$.
}
\end{ex}

% ===== Câu 56 =====
% ID: AIQ_L10C5_FULL_0218
% Mức: NB
\begin{ex}
% ID: AIQ_L10C5_FULL_0218
% Mức: NB
Vật $m_1=2$ kg chuyển động $5$ m/s va chạm mềm cùng chiều với vật $m_2=3$ kg chuyển động $1$ m/s. Tính vận tốc chung.
\choice
{$6\ \mathrm{m/s}$}
{\True $2,6\ \mathrm{m/s}$}
{$4\ \mathrm{m/s}$}
{$13\ \mathrm{m/s}$}
\loigiai{
$v'=\dfrac{m_1v_1+m_2v_2}{m_1+m_2}=\dfrac{2\cdot5+3\cdot1}{2+3}$. Suy ra kết quả $2,6\ \mathrm{m/s}$.
}
\end{ex}

% ===== Câu 57 =====
% ID: AIQ_L10C5_FULL_0219
% Mức: TH
\begin{ex}
% ID: AIQ_L10C5_FULL_0219
% Mức: TH
Vật $m_1=3$ kg chuyển động $6$ m/s va chạm mềm cùng chiều với vật $m_2=4$ kg chuyển động $0$ m/s. Tính vận tốc chung.
\choice
{$6\ \mathrm{m/s}$}
{$6\ \mathrm{m/s}$}
{\True $2,57\ \mathrm{m/s}$}
{$18\ \mathrm{m/s}$}
\loigiai{
$v'=\dfrac{m_1v_1+m_2v_2}{m_1+m_2}=\dfrac{3\cdot6+4\cdot0}{3+4}$. Suy ra kết quả $2,57\ \mathrm{m/s}$.
}
\end{ex}

% ===== Câu 58 =====
% ID: AIQ_L10C5_FULL_0220
% Mức: TH
\begin{ex}
% ID: AIQ_L10C5_FULL_0220
% Mức: TH
Vật $m_1=4$ kg chuyển động $7$ m/s va chạm mềm cùng chiều với vật $m_2=2$ kg chuyển động $1$ m/s. Tính vận tốc chung.
\choice
{$8\ \mathrm{m/s}$}
{$6\ \mathrm{m/s}$}
{$30\ \mathrm{m/s}$}
{\True $5\ \mathrm{m/s}$}
\loigiai{
$v'=\dfrac{m_1v_1+m_2v_2}{m_1+m_2}=\dfrac{4\cdot7+2\cdot1}{4+2}$. Suy ra kết quả $5\ \mathrm{m/s}$.
}
\end{ex}

% ===== Câu 59 =====
% ID: AIQ_L10C5_FULL_0221
% Mức: VD
\begin{ex}
% ID: AIQ_L10C5_FULL_0221
% Mức: VD
Vật $m_1=1$ kg chuyển động $8$ m/s va chạm mềm cùng chiều với vật $m_2=3$ kg chuyển động $0$ m/s. Tính vận tốc chung.
\choice
{\True $2\ \mathrm{m/s}$}
{$8\ \mathrm{m/s}$}
{$8\ \mathrm{m/s}$}
{$8\ \mathrm{m/s}$}
\loigiai{
$v'=\dfrac{m_1v_1+m_2v_2}{m_1+m_2}=\dfrac{1\cdot8+3\cdot0}{1+3}$. Suy ra kết quả $2\ \mathrm{m/s}$.
}
\end{ex}

% ===== Câu 60 =====
% ID: AIQ_L10C5_FULL_0222
% Mức: VD
\begin{ex}
% ID: AIQ_L10C5_FULL_0222
% Mức: VD
Vật $m_1=2$ kg chuyển động $9$ m/s va chạm mềm cùng chiều với vật $m_2=4$ kg chuyển động $1$ m/s. Tính vận tốc chung.
\choice
{$10\ \mathrm{m/s}$}
{\True $3,67\ \mathrm{m/s}$}
{$8\ \mathrm{m/s}$}
{$22\ \mathrm{m/s}$}
\loigiai{
$v'=\dfrac{m_1v_1+m_2v_2}{m_1+m_2}=\dfrac{2\cdot9+4\cdot1}{2+4}$. Suy ra kết quả $3,67\ \mathrm{m/s}$.
}
\end{ex}

% ===== Câu 61 =====
% ID: AIQ_L10C5_FULL_0223
% Mức: VD
\begin{ex}
% ID: AIQ_L10C5_FULL_0223
% Mức: VD
Vật $m_1=3$ kg chuyển động $10$ m/s va chạm mềm cùng chiều với vật $m_2=2$ kg chuyển động $0$ m/s. Tính vận tốc chung.
\choice
{$10\ \mathrm{m/s}$}
{$10\ \mathrm{m/s}$}
{\True $6\ \mathrm{m/s}$}
{$30\ \mathrm{m/s}$}
\loigiai{
$v'=\dfrac{m_1v_1+m_2v_2}{m_1+m_2}=\dfrac{3\cdot10+2\cdot0}{3+2}$. Suy ra kết quả $6\ \mathrm{m/s}$.
}
\end{ex}

% ===== Câu 62 =====
% ID: AIQ_L10C5_FULL_0224
% Mức: VDC
\begin{ex}
% ID: AIQ_L10C5_FULL_0224
% Mức: VDC
Vật $m_1=4$ kg chuyển động $11$ m/s va chạm mềm cùng chiều với vật $m_2=3$ kg chuyển động $1$ m/s. Tính vận tốc chung.
\choice
{$12\ \mathrm{m/s}$}
{$10\ \mathrm{m/s}$}
{$47\ \mathrm{m/s}$}
{\True $6,71\ \mathrm{m/s}$}
\loigiai{
$v'=\dfrac{m_1v_1+m_2v_2}{m_1+m_2}=\dfrac{4\cdot11+3\cdot1}{4+3}$. Suy ra kết quả $6,71\ \mathrm{m/s}$.
}
\end{ex}

% ===== Câu 63 =====
% ID: AIQ_L10C5_FULL_0225
% Mức: VDC
\begin{ex}
% ID: AIQ_L10C5_FULL_0225
% Mức: VDC
Vật $m_1=1$ kg chuyển động $12$ m/s va chạm mềm cùng chiều với vật $m_2=4$ kg chuyển động $0$ m/s. Tính vận tốc chung.
\choice
{\True $2,4\ \mathrm{m/s}$}
{$12\ \mathrm{m/s}$}
{$12\ \mathrm{m/s}$}
{$12\ \mathrm{m/s}$}
\loigiai{
$v'=\dfrac{m_1v_1+m_2v_2}{m_1+m_2}=\dfrac{1\cdot12+4\cdot0}{1+4}$. Suy ra kết quả $2,4\ \mathrm{m/s}$.
}
\end{ex}

% ===== Câu 64 =====
% ID: AIQ_L10C5_FULL_0226
% Mức: VDC
\begin{ex}
% ID: AIQ_L10C5_FULL_0226
% Mức: VDC
Vật $m_1=2$ kg chuyển động $13$ m/s va chạm mềm cùng chiều với vật $m_2=2$ kg chuyển động $1$ m/s. Tính vận tốc chung.
\choice
{$14\ \mathrm{m/s}$}
{\True $7\ \mathrm{m/s}$}
{$12\ \mathrm{m/s}$}
{$28\ \mathrm{m/s}$}
\loigiai{
$v'=\dfrac{m_1v_1+m_2v_2}{m_1+m_2}=\dfrac{2\cdot13+2\cdot1}{2+2}$. Suy ra kết quả $7\ \mathrm{m/s}$.
}
\end{ex}

% ===== Câu 65 =====
% ID: AIQ_L10C5_FULL_0227
% Mức: NB
\begin{ex}
% ID: AIQ_L10C5_FULL_0227
% Mức: NB
Xét các phát biểu thuộc dạng “Bảo toàn động lượng trong va chạm mềm” ở tình huống số 1.
\choiceTF
{\True Va chạm mềm làm các vật dính vào nhau.}
{Động năng luôn bảo toàn trong va chạm mềm.}
{\True Động lượng hệ kín vẫn bảo toàn.}
{\True Vận tốc chung được tính từ tổng động lượng chia tổng khối lượng.}
\loigiai{
\begin{itemchoice}
\item (Đúng) Va chạm mềm làm các vật dính vào nhau. (Vì): Phù hợp với định nghĩa hoặc hệ thức vật lí. b) Sai: Không phù hợp với định nghĩa hoặc hệ thức vật lí. c) Đúng: Phù hợp với định nghĩa hoặc hệ thức vật lí. d) Đúng: Phù hợp với định nghĩa hoặc hệ thức vật lí
\item (Sai) Động năng luôn bảo toàn trong va chạm mềm.
\item (Đúng) Động lượng hệ kín vẫn bảo toàn.
\item (Đúng) Vận tốc chung được tính từ tổng động lượng chia tổng khối lượng.
\end{itemchoice}
}
\end{ex}

% ===== Câu 66 =====
% ID: AIQ_L10C5_FULL_0228
% Mức: TH
\begin{ex}
% ID: AIQ_L10C5_FULL_0228
% Mức: TH
Xét các phát biểu thuộc dạng “Bảo toàn động lượng trong va chạm mềm” ở tình huống số 2.
\choiceTF
{\True Va chạm mềm làm các vật dính vào nhau.}
{Động năng luôn bảo toàn trong va chạm mềm.}
{\True Động lượng hệ kín vẫn bảo toàn.}
{\True Vận tốc chung được tính từ tổng động lượng chia tổng khối lượng.}
\loigiai{
\begin{itemchoice}
\item (Đúng) Va chạm mềm làm các vật dính vào nhau. (Vì): Phù hợp với định nghĩa hoặc hệ thức vật lí. b) Sai: Không phù hợp với định nghĩa hoặc hệ thức vật lí. c) Đúng: Phù hợp với định nghĩa hoặc hệ thức vật lí. d) Đúng: Phù hợp với định nghĩa hoặc hệ thức vật lí
\item (Sai) Động năng luôn bảo toàn trong va chạm mềm.
\item (Đúng) Động lượng hệ kín vẫn bảo toàn.
\item (Đúng) Vận tốc chung được tính từ tổng động lượng chia tổng khối lượng.
\end{itemchoice}
}
\end{ex}

% ===== Câu 67 =====
% ID: AIQ_L10C5_FULL_0229
% Mức: TH
\begin{ex}
% ID: AIQ_L10C5_FULL_0229
% Mức: TH
Xét các phát biểu thuộc dạng “Bảo toàn động lượng trong va chạm mềm” ở tình huống số 3.
\choiceTF
{\True Va chạm mềm làm các vật dính vào nhau.}
{Động năng luôn bảo toàn trong va chạm mềm.}
{\True Động lượng hệ kín vẫn bảo toàn.}
{\True Vận tốc chung được tính từ tổng động lượng chia tổng khối lượng.}
\loigiai{
\begin{itemchoice}
\item (Đúng) Va chạm mềm làm các vật dính vào nhau. (Vì): Phù hợp với định nghĩa hoặc hệ thức vật lí. b) Sai: Không phù hợp với định nghĩa hoặc hệ thức vật lí. c) Đúng: Phù hợp với định nghĩa hoặc hệ thức vật lí. d) Đúng: Phù hợp với định nghĩa hoặc hệ thức vật lí
\item (Sai) Động năng luôn bảo toàn trong va chạm mềm.
\item (Đúng) Động lượng hệ kín vẫn bảo toàn.
\item (Đúng) Vận tốc chung được tính từ tổng động lượng chia tổng khối lượng.
\end{itemchoice}
}
\end{ex}

% ===== Câu 68 =====
% ID: AIQ_L10C5_FULL_0230
% Mức: VD
\begin{ex}
% ID: AIQ_L10C5_FULL_0230
% Mức: VD
Xét các phát biểu thuộc dạng “Bảo toàn động lượng trong va chạm mềm” ở tình huống số 4.
\choiceTF
{\True Va chạm mềm làm các vật dính vào nhau.}
{Động năng luôn bảo toàn trong va chạm mềm.}
{\True Động lượng hệ kín vẫn bảo toàn.}
{\True Vận tốc chung được tính từ tổng động lượng chia tổng khối lượng.}
\loigiai{
\begin{itemchoice}
\item (Đúng) Va chạm mềm làm các vật dính vào nhau. (Vì): Phù hợp với định nghĩa hoặc hệ thức vật lí. b) Sai: Không phù hợp với định nghĩa hoặc hệ thức vật lí. c) Đúng: Phù hợp với định nghĩa hoặc hệ thức vật lí. d) Đúng: Phù hợp với định nghĩa hoặc hệ thức vật lí
\item (Sai) Động năng luôn bảo toàn trong va chạm mềm.
\item (Đúng) Động lượng hệ kín vẫn bảo toàn.
\item (Đúng) Vận tốc chung được tính từ tổng động lượng chia tổng khối lượng.
\end{itemchoice}
}
\end{ex}

% ===== Câu 69 =====
% ID: AIQ_L10C5_FULL_0231
% Mức: VDC
\begin{ex}
% ID: AIQ_L10C5_FULL_0231
% Mức: VDC
Xét các phát biểu thuộc dạng “Bảo toàn động lượng trong va chạm mềm” ở tình huống số 5.
\choiceTF
{\True Va chạm mềm làm các vật dính vào nhau.}
{Động năng luôn bảo toàn trong va chạm mềm.}
{\True Động lượng hệ kín vẫn bảo toàn.}
{\True Vận tốc chung được tính từ tổng động lượng chia tổng khối lượng.}
\loigiai{
\begin{itemchoice}
\item (Đúng) Va chạm mềm làm các vật dính vào nhau. (Vì): Phù hợp với định nghĩa hoặc hệ thức vật lí. b) Sai: Không phù hợp với định nghĩa hoặc hệ thức vật lí. c) Đúng: Phù hợp với định nghĩa hoặc hệ thức vật lí. d) Đúng: Phù hợp với định nghĩa hoặc hệ thức vật lí
\item (Sai) Động năng luôn bảo toàn trong va chạm mềm.
\item (Đúng) Động lượng hệ kín vẫn bảo toàn.
\item (Đúng) Vận tốc chung được tính từ tổng động lượng chia tổng khối lượng.
\end{itemchoice}
}
\end{ex}

% ===== Câu 70 =====
% ID: AIQ_L10C5_FULL_0232
% Mức: TH
\begin{ex}
% ID: AIQ_L10C5_FULL_0232
% Mức: TH
Vật $m_1=3$ kg chuyển động $14$ m/s va chạm mềm cùng chiều với vật $m_2=3$ kg chuyển động $0$ m/s. Tính vận tốc chung.
\shortans{7}
\loigiai{
$v'=\dfrac{m_1v_1+m_2v_2}{m_1+m_2}=\dfrac{3\cdot14+3\cdot0}{3+3}$. Suy ra kết quả $7\ \mathrm{m/s}$. Đáp án trả lời ngắn không ghi đơn vị.
}
\end{ex}

% ===== Câu 71 =====
% ID: AIQ_L10C5_FULL_0233
% Mức: TH
\begin{ex}
% ID: AIQ_L10C5_FULL_0233
% Mức: TH
Vật $m_1=4$ kg chuyển động $15$ m/s va chạm mềm cùng chiều với vật $m_2=4$ kg chuyển động $1$ m/s. Tính vận tốc chung.
\shortans{8}
\loigiai{
$v'=\dfrac{m_1v_1+m_2v_2}{m_1+m_2}=\dfrac{4\cdot15+4\cdot1}{4+4}$. Suy ra kết quả $8\ \mathrm{m/s}$. Đáp án trả lời ngắn không ghi đơn vị.
}
\end{ex}

% ===== Câu 72 =====
% ID: AIQ_L10C5_FULL_0234
% Mức: VD
\begin{ex}
% ID: AIQ_L10C5_FULL_0234
% Mức: VD
Vật $m_1=1$ kg chuyển động $16$ m/s va chạm mềm cùng chiều với vật $m_2=2$ kg chuyển động $0$ m/s. Tính vận tốc chung.
\shortans{5,33}
\loigiai{
$v'=\dfrac{m_1v_1+m_2v_2}{m_1+m_2}=\dfrac{1\cdot16+2\cdot0}{1+2}$. Suy ra kết quả $5,33\ \mathrm{m/s}$. Đáp án trả lời ngắn không ghi đơn vị.
}
\end{ex}

% ===== Câu 73 =====
% ID: AIQ_L10C5_FULL_0235
% Mức: VD
\begin{ex}
% ID: AIQ_L10C5_FULL_0235
% Mức: VD
Vật $m_1=2$ kg chuyển động $17$ m/s va chạm mềm cùng chiều với vật $m_2=3$ kg chuyển động $1$ m/s. Tính vận tốc chung.
\shortans{7,4}
\loigiai{
$v'=\dfrac{m_1v_1+m_2v_2}{m_1+m_2}=\dfrac{2\cdot17+3\cdot1}{2+3}$. Suy ra kết quả $7,4\ \mathrm{m/s}$. Đáp án trả lời ngắn không ghi đơn vị.
}
\end{ex}

% ===== Câu 74 =====
% ID: AIQ_L10C5_FULL_0236
% Mức: VDC
\begin{ex}
% ID: AIQ_L10C5_FULL_0236
% Mức: VDC
Vật $m_1=3$ kg chuyển động $18$ m/s va chạm mềm cùng chiều với vật $m_2=4$ kg chuyển động $0$ m/s. Tính vận tốc chung.
\shortans{7,71}
\loigiai{
$v'=\dfrac{m_1v_1+m_2v_2}{m_1+m_2}=\dfrac{3\cdot18+4\cdot0}{3+4}$. Suy ra kết quả $7,71\ \mathrm{m/s}$. Đáp án trả lời ngắn không ghi đơn vị.
}
\end{ex}

% ===== Câu 75 =====
% ID: AIQ_L10C5_FULL_0237
% Mức: VDC
\begin{ex}
% ID: AIQ_L10C5_FULL_0237
% Mức: VDC
Vật $m_1=4$ kg chuyển động $19$ m/s va chạm mềm cùng chiều với vật $m_2=2$ kg chuyển động $1$ m/s. Tính vận tốc chung.
\shortans{13}
\loigiai{
$v'=\dfrac{m_1v_1+m_2v_2}{m_1+m_2}=\dfrac{4\cdot19+2\cdot1}{4+2}$. Suy ra kết quả $13\ \mathrm{m/s}$. Đáp án trả lời ngắn không ghi đơn vị.
}
\end{ex}

% ===== Câu 76 =====
% ID: AIQ_L10C5_FULL_0238
% Mức: TH
\begin{bt}
% ID: AIQ_L10C5_FULL_0238
% Mức: TH
Trình bày cách giải: Vật $m_1=1$ kg chuyển động $20$ m/s va chạm mềm cùng chiều với vật $m_2=3$ kg chuyển động $0$ m/s. Tính vận tốc chung.
\loigiai{
$v'=\dfrac{m_1v_1+m_2v_2}{m_1+m_2}=\dfrac{1\cdot20+3\cdot0}{1+3}$. Suy ra kết quả $5\ \mathrm{m/s}$.
}
\end{bt}

% ===== Câu 77 =====
% ID: AIQ_L10C5_FULL_0239
% Mức: TH
\begin{bt}
% ID: AIQ_L10C5_FULL_0239
% Mức: TH
Trình bày cách giải: Vật $m_1=2$ kg chuyển động $21$ m/s va chạm mềm cùng chiều với vật $m_2=4$ kg chuyển động $1$ m/s. Tính vận tốc chung.
\loigiai{
$v'=\dfrac{m_1v_1+m_2v_2}{m_1+m_2}=\dfrac{2\cdot21+4\cdot1}{2+4}$. Suy ra kết quả $7,67\ \mathrm{m/s}$.
}
\end{bt}

% ===== Câu 78 =====
% ID: AIQ_L10C5_FULL_0240
% Mức: VD
\begin{bt}
% ID: AIQ_L10C5_FULL_0240
% Mức: VD
Trình bày cách giải: Vật $m_1=3$ kg chuyển động $22$ m/s va chạm mềm cùng chiều với vật $m_2=2$ kg chuyển động $0$ m/s. Tính vận tốc chung.
\loigiai{
$v'=\dfrac{m_1v_1+m_2v_2}{m_1+m_2}=\dfrac{3\cdot22+2\cdot0}{3+2}$. Suy ra kết quả $13,2\ \mathrm{m/s}$.
}
\end{bt}

% ===== Câu 79 =====
% ID: AIQ_L10C5_FULL_0241
% Mức: VD
\begin{bt}
% ID: AIQ_L10C5_FULL_0241
% Mức: VD
Trình bày cách giải: Vật $m_1=4$ kg chuyển động $23$ m/s va chạm mềm cùng chiều với vật $m_2=3$ kg chuyển động $1$ m/s. Tính vận tốc chung.
\loigiai{
$v'=\dfrac{m_1v_1+m_2v_2}{m_1+m_2}=\dfrac{4\cdot23+3\cdot1}{4+3}$. Suy ra kết quả $13,57\ \mathrm{m/s}$.
}
\end{bt}

% ===== Câu 80 =====
% ID: AIQ_L10C5_FULL_0242
% Mức: VDC
\begin{bt}
% ID: AIQ_L10C5_FULL_0242
% Mức: VDC
Trình bày cách giải: Vật $m_1=1$ kg chuyển động $24$ m/s va chạm mềm cùng chiều với vật $m_2=4$ kg chuyển động $0$ m/s. Tính vận tốc chung.
\loigiai{
$v'=\dfrac{m_1v_1+m_2v_2}{m_1+m_2}=\dfrac{1\cdot24+4\cdot0}{1+4}$. Suy ra kết quả $4,8\ \mathrm{m/s}$.
}
\end{bt}

% ===== Câu 81 =====
% ID: AIQ_L10C5_FULL_0243
% Mức: VDC
\begin{bt}
% ID: AIQ_L10C5_FULL_0243
% Mức: VDC
Trình bày cách giải: Vật $m_1=2$ kg chuyển động $25$ m/s va chạm mềm cùng chiều với vật $m_2=2$ kg chuyển động $1$ m/s. Tính vận tốc chung.
\loigiai{
$v'=\dfrac{m_1v_1+m_2v_2}{m_1+m_2}=\dfrac{2\cdot25+2\cdot1}{2+2}$. Suy ra kết quả $13\ \mathrm{m/s}$.
}
\end{bt}

% ===== Câu 82 =====
% ID: AIQ_L10C5_FULL_0244
% Mức: NB
\dangbt{Bảo toàn động lượng trong chuyển động phản lực và nổ}
\begin{ex}
% ID: AIQ_L10C5_FULL_0244
% Mức: NB
Một vật khối lượng $1$ kg được phóng theo chiều dương với tốc độ $6$ m/s từ xe khối lượng $4$ kg ban đầu đứng yên. Tính độ lớn vận tốc giật lùi của xe.
\choice
{\True $1,5\ \mathrm{m/s}$}
{$6\ \mathrm{m/s}$}
{$24\ \mathrm{m/s}$}
{$3\ \mathrm{m/s}$}
\loigiai{
Bảo toàn động lượng: $mv-MV=0\Rightarrow V=mv/M=1\cdot6/4$. Suy ra kết quả $1,5\ \mathrm{m/s}$.
}
\end{ex}

% ===== Câu 83 =====
% ID: AIQ_L10C5_FULL_0245
% Mức: NB
\begin{ex}
% ID: AIQ_L10C5_FULL_0245
% Mức: NB
Một vật khối lượng $2$ kg được phóng theo chiều dương với tốc độ $7$ m/s từ xe khối lượng $5$ kg ban đầu đứng yên. Tính độ lớn vận tốc giật lùi của xe.
\choice
{$14\ \mathrm{m/s}$}
{\True $2,8\ \mathrm{m/s}$}
{$35\ \mathrm{m/s}$}
{$5,6\ \mathrm{m/s}$}
\loigiai{
Bảo toàn động lượng: $mv-MV=0\Rightarrow V=mv/M=2\cdot7/5$. Suy ra kết quả $2,8\ \mathrm{m/s}$.
}
\end{ex}

% ===== Câu 84 =====
% ID: AIQ_L10C5_FULL_0246
% Mức: TH
\begin{ex}
% ID: AIQ_L10C5_FULL_0246
% Mức: TH
Một vật khối lượng $3$ kg được phóng theo chiều dương với tốc độ $8$ m/s từ xe khối lượng $6$ kg ban đầu đứng yên. Tính độ lớn vận tốc giật lùi của xe.
\choice
{$24\ \mathrm{m/s}$}
{$48\ \mathrm{m/s}$}
{\True $4\ \mathrm{m/s}$}
{$8\ \mathrm{m/s}$}
\loigiai{
Bảo toàn động lượng: $mv-MV=0\Rightarrow V=mv/M=3\cdot8/6$. Suy ra kết quả $4\ \mathrm{m/s}$.
}
\end{ex}

% ===== Câu 85 =====
% ID: AIQ_L10C5_FULL_0247
% Mức: TH
\begin{ex}
% ID: AIQ_L10C5_FULL_0247
% Mức: TH
Một vật khối lượng $1$ kg được phóng theo chiều dương với tốc độ $9$ m/s từ xe khối lượng $7$ kg ban đầu đứng yên. Tính độ lớn vận tốc giật lùi của xe.
\choice
{$9\ \mathrm{m/s}$}
{$63\ \mathrm{m/s}$}
{$2,57\ \mathrm{m/s}$}
{\True $1,29\ \mathrm{m/s}$}
\loigiai{
Bảo toàn động lượng: $mv-MV=0\Rightarrow V=mv/M=1\cdot9/7$. Suy ra kết quả $1,29\ \mathrm{m/s}$.
}
\end{ex}

% ===== Câu 86 =====
% ID: AIQ_L10C5_FULL_0248
% Mức: VD
\begin{ex}
% ID: AIQ_L10C5_FULL_0248
% Mức: VD
Một vật khối lượng $2$ kg được phóng theo chiều dương với tốc độ $10$ m/s từ xe khối lượng $8$ kg ban đầu đứng yên. Tính độ lớn vận tốc giật lùi của xe.
\choice
{\True $2,5\ \mathrm{m/s}$}
{$20\ \mathrm{m/s}$}
{$80\ \mathrm{m/s}$}
{$5\ \mathrm{m/s}$}
\loigiai{
Bảo toàn động lượng: $mv-MV=0\Rightarrow V=mv/M=2\cdot10/8$. Suy ra kết quả $2,5\ \mathrm{m/s}$.
}
\end{ex}

% ===== Câu 87 =====
% ID: AIQ_L10C5_FULL_0249
% Mức: VD
\begin{ex}
% ID: AIQ_L10C5_FULL_0249
% Mức: VD
Một vật khối lượng $3$ kg được phóng theo chiều dương với tốc độ $11$ m/s từ xe khối lượng $4$ kg ban đầu đứng yên. Tính độ lớn vận tốc giật lùi của xe.
\choice
{$33\ \mathrm{m/s}$}
{\True $8,25\ \mathrm{m/s}$}
{$44\ \mathrm{m/s}$}
{$16,5\ \mathrm{m/s}$}
\loigiai{
Bảo toàn động lượng: $mv-MV=0\Rightarrow V=mv/M=3\cdot11/4$. Suy ra kết quả $8,25\ \mathrm{m/s}$.
}
\end{ex}

% ===== Câu 88 =====
% ID: AIQ_L10C5_FULL_0250
% Mức: VD
\begin{ex}
% ID: AIQ_L10C5_FULL_0250
% Mức: VD
Một vật khối lượng $1$ kg được phóng theo chiều dương với tốc độ $12$ m/s từ xe khối lượng $5$ kg ban đầu đứng yên. Tính độ lớn vận tốc giật lùi của xe.
\choice
{$12\ \mathrm{m/s}$}
{$60\ \mathrm{m/s}$}
{\True $2,4\ \mathrm{m/s}$}
{$4,8\ \mathrm{m/s}$}
\loigiai{
Bảo toàn động lượng: $mv-MV=0\Rightarrow V=mv/M=1\cdot12/5$. Suy ra kết quả $2,4\ \mathrm{m/s}$.
}
\end{ex}

% ===== Câu 89 =====
% ID: AIQ_L10C5_FULL_0251
% Mức: VDC
\begin{ex}
% ID: AIQ_L10C5_FULL_0251
% Mức: VDC
Một vật khối lượng $2$ kg được phóng theo chiều dương với tốc độ $13$ m/s từ xe khối lượng $6$ kg ban đầu đứng yên. Tính độ lớn vận tốc giật lùi của xe.
\choice
{$26\ \mathrm{m/s}$}
{$78\ \mathrm{m/s}$}
{$8,67\ \mathrm{m/s}$}
{\True $4,33\ \mathrm{m/s}$}
\loigiai{
Bảo toàn động lượng: $mv-MV=0\Rightarrow V=mv/M=2\cdot13/6$. Suy ra kết quả $4,33\ \mathrm{m/s}$.
}
\end{ex}

% ===== Câu 90 =====
% ID: AIQ_L10C5_FULL_0252
% Mức: VDC
\begin{ex}
% ID: AIQ_L10C5_FULL_0252
% Mức: VDC
Một vật khối lượng $3$ kg được phóng theo chiều dương với tốc độ $14$ m/s từ xe khối lượng $7$ kg ban đầu đứng yên. Tính độ lớn vận tốc giật lùi của xe.
\choice
{\True $6\ \mathrm{m/s}$}
{$42\ \mathrm{m/s}$}
{$98\ \mathrm{m/s}$}
{$12\ \mathrm{m/s}$}
\loigiai{
Bảo toàn động lượng: $mv-MV=0\Rightarrow V=mv/M=3\cdot14/7$. Suy ra kết quả $6\ \mathrm{m/s}$.
}
\end{ex}

% ===== Câu 91 =====
% ID: AIQ_L10C5_FULL_0253
% Mức: VDC
\begin{ex}
% ID: AIQ_L10C5_FULL_0253
% Mức: VDC
Một vật khối lượng $1$ kg được phóng theo chiều dương với tốc độ $15$ m/s từ xe khối lượng $8$ kg ban đầu đứng yên. Tính độ lớn vận tốc giật lùi của xe.
\choice
{$15\ \mathrm{m/s}$}
{\True $1,88\ \mathrm{m/s}$}
{$120\ \mathrm{m/s}$}
{$3,75\ \mathrm{m/s}$}
\loigiai{
Bảo toàn động lượng: $mv-MV=0\Rightarrow V=mv/M=1\cdot15/8$. Suy ra kết quả $1,88\ \mathrm{m/s}$.
}
\end{ex}

% ===== Câu 92 =====
% ID: AIQ_L10C5_FULL_0254
% Mức: NB
\begin{ex}
% ID: AIQ_L10C5_FULL_0254
% Mức: NB
Xét các phát biểu thuộc dạng “Bảo toàn động lượng trong chuyển động phản lực và nổ” ở tình huống số 1.
\choiceTF
{\True Chuyển động phản lực tuân theo bảo toàn động lượng.}
{\True Súng giật lùi khi bắn đạn.}
{Hai mảnh nổ luôn có cùng tốc độ.}
{\True Hệ ban đầu đứng yên có tổng động lượng sau nổ bằng 0.}
\loigiai{
\begin{itemchoice}
\item (Đúng) Chuyển động phản lực tuân theo bảo toàn động lượng. (Vì): Phù hợp với định nghĩa hoặc hệ thức vật lí. b) Đúng: Phù hợp với định nghĩa hoặc hệ thức vật lí. c) Sai: Không phù hợp với định nghĩa hoặc hệ thức vật lí. d) Đúng: Phù hợp với định nghĩa hoặc hệ thức vật lí
\item (Đúng) Súng giật lùi khi bắn đạn.
\item (Sai) Hai mảnh nổ luôn có cùng tốc độ.
\item (Đúng) Hệ ban đầu đứng yên có tổng động lượng sau nổ bằng 0.
\end{itemchoice}
}
\end{ex}

% ===== Câu 93 =====
% ID: AIQ_L10C5_FULL_0255
% Mức: TH
\begin{ex}
% ID: AIQ_L10C5_FULL_0255
% Mức: TH
Xét các phát biểu thuộc dạng “Bảo toàn động lượng trong chuyển động phản lực và nổ” ở tình huống số 2.
\choiceTF
{\True Chuyển động phản lực tuân theo bảo toàn động lượng.}
{\True Súng giật lùi khi bắn đạn.}
{Hai mảnh nổ luôn có cùng tốc độ.}
{\True Hệ ban đầu đứng yên có tổng động lượng sau nổ bằng 0.}
\loigiai{
\begin{itemchoice}
\item (Đúng) Chuyển động phản lực tuân theo bảo toàn động lượng. (Vì): Phù hợp với định nghĩa hoặc hệ thức vật lí. b) Đúng: Phù hợp với định nghĩa hoặc hệ thức vật lí. c) Sai: Không phù hợp với định nghĩa hoặc hệ thức vật lí. d) Đúng: Phù hợp với định nghĩa hoặc hệ thức vật lí
\item (Đúng) Súng giật lùi khi bắn đạn.
\item (Sai) Hai mảnh nổ luôn có cùng tốc độ.
\item (Đúng) Hệ ban đầu đứng yên có tổng động lượng sau nổ bằng 0.
\end{itemchoice}
}
\end{ex}

% ===== Câu 94 =====
% ID: AIQ_L10C5_FULL_0256
% Mức: TH
\begin{ex}
% ID: AIQ_L10C5_FULL_0256
% Mức: TH
Xét các phát biểu thuộc dạng “Bảo toàn động lượng trong chuyển động phản lực và nổ” ở tình huống số 3.
\choiceTF
{\True Chuyển động phản lực tuân theo bảo toàn động lượng.}
{\True Súng giật lùi khi bắn đạn.}
{Hai mảnh nổ luôn có cùng tốc độ.}
{\True Hệ ban đầu đứng yên có tổng động lượng sau nổ bằng 0.}
\loigiai{
\begin{itemchoice}
\item (Đúng) Chuyển động phản lực tuân theo bảo toàn động lượng. (Vì): Phù hợp với định nghĩa hoặc hệ thức vật lí. b) Đúng: Phù hợp với định nghĩa hoặc hệ thức vật lí. c) Sai: Không phù hợp với định nghĩa hoặc hệ thức vật lí. d) Đúng: Phù hợp với định nghĩa hoặc hệ thức vật lí
\item (Đúng) Súng giật lùi khi bắn đạn.
\item (Sai) Hai mảnh nổ luôn có cùng tốc độ.
\item (Đúng) Hệ ban đầu đứng yên có tổng động lượng sau nổ bằng 0.
\end{itemchoice}
}
\end{ex}

% ===== Câu 95 =====
% ID: AIQ_L10C5_FULL_0257
% Mức: VD
\begin{ex}
% ID: AIQ_L10C5_FULL_0257
% Mức: VD
Xét các phát biểu thuộc dạng “Bảo toàn động lượng trong chuyển động phản lực và nổ” ở tình huống số 4.
\choiceTF
{\True Chuyển động phản lực tuân theo bảo toàn động lượng.}
{\True Súng giật lùi khi bắn đạn.}
{Hai mảnh nổ luôn có cùng tốc độ.}
{\True Hệ ban đầu đứng yên có tổng động lượng sau nổ bằng 0.}
\loigiai{
\begin{itemchoice}
\item (Đúng) Chuyển động phản lực tuân theo bảo toàn động lượng. (Vì): Phù hợp với định nghĩa hoặc hệ thức vật lí. b) Đúng: Phù hợp với định nghĩa hoặc hệ thức vật lí. c) Sai: Không phù hợp với định nghĩa hoặc hệ thức vật lí. d) Đúng: Phù hợp với định nghĩa hoặc hệ thức vật lí
\item (Đúng) Súng giật lùi khi bắn đạn.
\item (Sai) Hai mảnh nổ luôn có cùng tốc độ.
\item (Đúng) Hệ ban đầu đứng yên có tổng động lượng sau nổ bằng 0.
\end{itemchoice}
}
\end{ex}

% ===== Câu 96 =====
% ID: AIQ_L10C5_FULL_0258
% Mức: VDC
\begin{ex}
% ID: AIQ_L10C5_FULL_0258
% Mức: VDC
Xét các phát biểu thuộc dạng “Bảo toàn động lượng trong chuyển động phản lực và nổ” ở tình huống số 5.
\choiceTF
{\True Chuyển động phản lực tuân theo bảo toàn động lượng.}
{\True Súng giật lùi khi bắn đạn.}
{Hai mảnh nổ luôn có cùng tốc độ.}
{\True Hệ ban đầu đứng yên có tổng động lượng sau nổ bằng 0.}
\loigiai{
\begin{itemchoice}
\item (Đúng) Chuyển động phản lực tuân theo bảo toàn động lượng. (Vì): Phù hợp với định nghĩa hoặc hệ thức vật lí. b) Đúng: Phù hợp với định nghĩa hoặc hệ thức vật lí. c) Sai: Không phù hợp với định nghĩa hoặc hệ thức vật lí. d) Đúng: Phù hợp với định nghĩa hoặc hệ thức vật lí
\item (Đúng) Súng giật lùi khi bắn đạn.
\item (Sai) Hai mảnh nổ luôn có cùng tốc độ.
\item (Đúng) Hệ ban đầu đứng yên có tổng động lượng sau nổ bằng 0.
\end{itemchoice}
}
\end{ex}

% ===== Câu 97 =====
% ID: AIQ_L10C5_FULL_0259
% Mức: TH
\begin{ex}
% ID: AIQ_L10C5_FULL_0259
% Mức: TH
Một vật khối lượng $2$ kg được phóng theo chiều dương với tốc độ $16$ m/s từ xe khối lượng $4$ kg ban đầu đứng yên. Tính độ lớn vận tốc giật lùi của xe.
\shortans{8}
\loigiai{
Bảo toàn động lượng: $mv-MV=0\Rightarrow V=mv/M=2\cdot16/4$. Suy ra kết quả $8\ \mathrm{m/s}$. Đáp án trả lời ngắn không ghi đơn vị.
}
\end{ex}

% ===== Câu 98 =====
% ID: AIQ_L10C5_FULL_0260
% Mức: TH
\begin{ex}
% ID: AIQ_L10C5_FULL_0260
% Mức: TH
Một vật khối lượng $3$ kg được phóng theo chiều dương với tốc độ $17$ m/s từ xe khối lượng $5$ kg ban đầu đứng yên. Tính độ lớn vận tốc giật lùi của xe.
\shortans{10,2}
\loigiai{
Bảo toàn động lượng: $mv-MV=0\Rightarrow V=mv/M=3\cdot17/5$. Suy ra kết quả $10,2\ \mathrm{m/s}$. Đáp án trả lời ngắn không ghi đơn vị.
}
\end{ex}

% ===== Câu 99 =====
% ID: AIQ_L10C5_FULL_0261
% Mức: VD
\begin{ex}
% ID: AIQ_L10C5_FULL_0261
% Mức: VD
Một vật khối lượng $1$ kg được phóng theo chiều dương với tốc độ $18$ m/s từ xe khối lượng $6$ kg ban đầu đứng yên. Tính độ lớn vận tốc giật lùi của xe.
\shortans{3}
\loigiai{
Bảo toàn động lượng: $mv-MV=0\Rightarrow V=mv/M=1\cdot18/6$. Suy ra kết quả $3\ \mathrm{m/s}$. Đáp án trả lời ngắn không ghi đơn vị.
}
\end{ex}

% ===== Câu 100 =====
% ID: AIQ_L10C5_FULL_0262
% Mức: VD
\begin{ex}
% ID: AIQ_L10C5_FULL_0262
% Mức: VD
Một vật khối lượng $2$ kg được phóng theo chiều dương với tốc độ $19$ m/s từ xe khối lượng $7$ kg ban đầu đứng yên. Tính độ lớn vận tốc giật lùi của xe.
\shortans{5,43}
\loigiai{
Bảo toàn động lượng: $mv-MV=0\Rightarrow V=mv/M=2\cdot19/7$. Suy ra kết quả $5,43\ \mathrm{m/s}$. Đáp án trả lời ngắn không ghi đơn vị.
}
\end{ex}

% ===== Câu 101 =====
% ID: AIQ_L10C5_FULL_0263
% Mức: VDC
\begin{ex}
% ID: AIQ_L10C5_FULL_0263
% Mức: VDC
Một vật khối lượng $3$ kg được phóng theo chiều dương với tốc độ $20$ m/s từ xe khối lượng $8$ kg ban đầu đứng yên. Tính độ lớn vận tốc giật lùi của xe.
\shortans{7,5}
\loigiai{
Bảo toàn động lượng: $mv-MV=0\Rightarrow V=mv/M=3\cdot20/8$. Suy ra kết quả $7,5\ \mathrm{m/s}$. Đáp án trả lời ngắn không ghi đơn vị.
}
\end{ex}

% ===== Câu 102 =====
% ID: AIQ_L10C5_FULL_0264
% Mức: VDC
\begin{ex}
% ID: AIQ_L10C5_FULL_0264
% Mức: VDC
Một vật khối lượng $1$ kg được phóng theo chiều dương với tốc độ $21$ m/s từ xe khối lượng $4$ kg ban đầu đứng yên. Tính độ lớn vận tốc giật lùi của xe.
\shortans{5,25}
\loigiai{
Bảo toàn động lượng: $mv-MV=0\Rightarrow V=mv/M=1\cdot21/4$. Suy ra kết quả $5,25\ \mathrm{m/s}$. Đáp án trả lời ngắn không ghi đơn vị.
}
\end{ex}

% ===== Câu 103 =====
% ID: AIQ_L10C5_FULL_0265
% Mức: TH
\begin{bt}
% ID: AIQ_L10C5_FULL_0265
% Mức: TH
Trình bày cách giải: Một vật khối lượng $2$ kg được phóng theo chiều dương với tốc độ $22$ m/s từ xe khối lượng $5$ kg ban đầu đứng yên. Tính độ lớn vận tốc giật lùi của xe.
\loigiai{
Bảo toàn động lượng: $mv-MV=0\Rightarrow V=mv/M=2\cdot22/5$. Suy ra kết quả $8,8\ \mathrm{m/s}$.
}
\end{bt}

% ===== Câu 104 =====
% ID: AIQ_L10C5_FULL_0266
% Mức: TH
\begin{bt}
% ID: AIQ_L10C5_FULL_0266
% Mức: TH
Trình bày cách giải: Một vật khối lượng $3$ kg được phóng theo chiều dương với tốc độ $23$ m/s từ xe khối lượng $6$ kg ban đầu đứng yên. Tính độ lớn vận tốc giật lùi của xe.
\loigiai{
Bảo toàn động lượng: $mv-MV=0\Rightarrow V=mv/M=3\cdot23/6$. Suy ra kết quả $11,5\ \mathrm{m/s}$.
}
\end{bt}

% ===== Câu 105 =====
% ID: AIQ_L10C5_FULL_0267
% Mức: VD
\begin{bt}
% ID: AIQ_L10C5_FULL_0267
% Mức: VD
Trình bày cách giải: Một vật khối lượng $1$ kg được phóng theo chiều dương với tốc độ $24$ m/s từ xe khối lượng $7$ kg ban đầu đứng yên. Tính độ lớn vận tốc giật lùi của xe.
\loigiai{
Bảo toàn động lượng: $mv-MV=0\Rightarrow V=mv/M=1\cdot24/7$. Suy ra kết quả $3,43\ \mathrm{m/s}$.
}
\end{bt}

% ===== Câu 106 =====
% ID: AIQ_L10C5_FULL_0268
% Mức: VD
\begin{bt}
% ID: AIQ_L10C5_FULL_0268
% Mức: VD
Trình bày cách giải: Một vật khối lượng $2$ kg được phóng theo chiều dương với tốc độ $25$ m/s từ xe khối lượng $8$ kg ban đầu đứng yên. Tính độ lớn vận tốc giật lùi của xe.
\loigiai{
Bảo toàn động lượng: $mv-MV=0\Rightarrow V=mv/M=2\cdot25/8$. Suy ra kết quả $6,25\ \mathrm{m/s}$.
}
\end{bt}

% ===== Câu 107 =====
% ID: AIQ_L10C5_FULL_0269
% Mức: VDC
\begin{bt}
% ID: AIQ_L10C5_FULL_0269
% Mức: VDC
Trình bày cách giải: Một vật khối lượng $3$ kg được phóng theo chiều dương với tốc độ $26$ m/s từ xe khối lượng $4$ kg ban đầu đứng yên. Tính độ lớn vận tốc giật lùi của xe.
\loigiai{
Bảo toàn động lượng: $mv-MV=0\Rightarrow V=mv/M=3\cdot26/4$. Suy ra kết quả $19,5\ \mathrm{m/s}$.
}
\end{bt}

% ===== Câu 108 =====
% ID: AIQ_L10C5_FULL_0270
% Mức: VDC
\begin{bt}
% ID: AIQ_L10C5_FULL_0270
% Mức: VDC
Trình bày cách giải: Một vật khối lượng $1$ kg được phóng theo chiều dương với tốc độ $27$ m/s từ xe khối lượng $5$ kg ban đầu đứng yên. Tính độ lớn vận tốc giật lùi của xe.
\loigiai{
Bảo toàn động lượng: $mv-MV=0\Rightarrow V=mv/M=1\cdot27/5$. Suy ra kết quả $5,4\ \mathrm{m/s}$.
}
\end{bt}

% ===== Câu 109 =====
% ID: AIQ_L10C5_FULL_0271
% Mức: NB
\dangbt{Bài toán động lượng hai chiều}
\begin{ex}
% ID: AIQ_L10C5_FULL_0271
% Mức: NB
Hai mảnh văng theo hai phương vuông góc có động lượng $3$ và $4\ \mathrm{kg\,m/s}$. Tính độ lớn tổng động lượng.
\choice
{\True $5\ \mathrm{kg\,m/s}$}
{$7\ \mathrm{kg\,m/s}$}
{$1\ \mathrm{kg\,m/s}$}
{$12\ \mathrm{kg\,m/s}$}
\loigiai{
Hai vectơ vuông góc nên $p=\sqrt{p_1^2+p_2^2}=\sqrt{3^2+4^2}$. Suy ra kết quả $5\ \mathrm{kg\,m/s}$.
}
\end{ex}

% ===== Câu 110 =====
% ID: AIQ_L10C5_FULL_0272
% Mức: NB
\begin{ex}
% ID: AIQ_L10C5_FULL_0272
% Mức: NB
Hai mảnh văng theo hai phương vuông góc có động lượng $4$ và $5\ \mathrm{kg\,m/s}$. Tính độ lớn tổng động lượng.
\choice
{$9\ \mathrm{kg\,m/s}$}
{\True $6,4\ \mathrm{kg\,m/s}$}
{$1\ \mathrm{kg\,m/s}$}
{$20\ \mathrm{kg\,m/s}$}
\loigiai{
Hai vectơ vuông góc nên $p=\sqrt{p_1^2+p_2^2}=\sqrt{4^2+5^2}$. Suy ra kết quả $6,4\ \mathrm{kg\,m/s}$.
}
\end{ex}

% ===== Câu 111 =====
% ID: AIQ_L10C5_FULL_0273
% Mức: TH
\begin{ex}
% ID: AIQ_L10C5_FULL_0273
% Mức: TH
Hai mảnh văng theo hai phương vuông góc có động lượng $5$ và $6\ \mathrm{kg\,m/s}$. Tính độ lớn tổng động lượng.
\choice
{$11\ \mathrm{kg\,m/s}$}
{$1\ \mathrm{kg\,m/s}$}
{\True $7,81\ \mathrm{kg\,m/s}$}
{$30\ \mathrm{kg\,m/s}$}
\loigiai{
Hai vectơ vuông góc nên $p=\sqrt{p_1^2+p_2^2}=\sqrt{5^2+6^2}$. Suy ra kết quả $7,81\ \mathrm{kg\,m/s}$.
}
\end{ex}

% ===== Câu 112 =====
% ID: AIQ_L10C5_FULL_0274
% Mức: TH
\begin{ex}
% ID: AIQ_L10C5_FULL_0274
% Mức: TH
Hai mảnh văng theo hai phương vuông góc có động lượng $6$ và $7\ \mathrm{kg\,m/s}$. Tính độ lớn tổng động lượng.
\choice
{$13\ \mathrm{kg\,m/s}$}
{$1\ \mathrm{kg\,m/s}$}
{$42\ \mathrm{kg\,m/s}$}
{\True $9,22\ \mathrm{kg\,m/s}$}
\loigiai{
Hai vectơ vuông góc nên $p=\sqrt{p_1^2+p_2^2}=\sqrt{6^2+7^2}$. Suy ra kết quả $9,22\ \mathrm{kg\,m/s}$.
}
\end{ex}

% ===== Câu 113 =====
% ID: AIQ_L10C5_FULL_0275
% Mức: VD
\begin{ex}
% ID: AIQ_L10C5_FULL_0275
% Mức: VD
Hai mảnh văng theo hai phương vuông góc có động lượng $7$ và $8\ \mathrm{kg\,m/s}$. Tính độ lớn tổng động lượng.
\choice
{\True $10,63\ \mathrm{kg\,m/s}$}
{$15\ \mathrm{kg\,m/s}$}
{$1\ \mathrm{kg\,m/s}$}
{$56\ \mathrm{kg\,m/s}$}
\loigiai{
Hai vectơ vuông góc nên $p=\sqrt{p_1^2+p_2^2}=\sqrt{7^2+8^2}$. Suy ra kết quả $10,63\ \mathrm{kg\,m/s}$.
}
\end{ex}

% ===== Câu 114 =====
% ID: AIQ_L10C5_FULL_0276
% Mức: VD
\begin{ex}
% ID: AIQ_L10C5_FULL_0276
% Mức: VD
Hai mảnh văng theo hai phương vuông góc có động lượng $8$ và $4\ \mathrm{kg\,m/s}$. Tính độ lớn tổng động lượng.
\choice
{$12\ \mathrm{kg\,m/s}$}
{\True $8,94\ \mathrm{kg\,m/s}$}
{$4\ \mathrm{kg\,m/s}$}
{$32\ \mathrm{kg\,m/s}$}
\loigiai{
Hai vectơ vuông góc nên $p=\sqrt{p_1^2+p_2^2}=\sqrt{8^2+4^2}$. Suy ra kết quả $8,94\ \mathrm{kg\,m/s}$.
}
\end{ex}

% ===== Câu 115 =====
% ID: AIQ_L10C5_FULL_0277
% Mức: VD
\begin{ex}
% ID: AIQ_L10C5_FULL_0277
% Mức: VD
Hai mảnh văng theo hai phương vuông góc có động lượng $9$ và $5\ \mathrm{kg\,m/s}$. Tính độ lớn tổng động lượng.
\choice
{$14\ \mathrm{kg\,m/s}$}
{$4\ \mathrm{kg\,m/s}$}
{\True $10,3\ \mathrm{kg\,m/s}$}
{$45\ \mathrm{kg\,m/s}$}
\loigiai{
Hai vectơ vuông góc nên $p=\sqrt{p_1^2+p_2^2}=\sqrt{9^2+5^2}$. Suy ra kết quả $10,3\ \mathrm{kg\,m/s}$.
}
\end{ex}

% ===== Câu 116 =====
% ID: AIQ_L10C5_FULL_0278
% Mức: VDC
\begin{ex}
% ID: AIQ_L10C5_FULL_0278
% Mức: VDC
Hai mảnh văng theo hai phương vuông góc có động lượng $10$ và $6\ \mathrm{kg\,m/s}$. Tính độ lớn tổng động lượng.
\choice
{$16\ \mathrm{kg\,m/s}$}
{$4\ \mathrm{kg\,m/s}$}
{$60\ \mathrm{kg\,m/s}$}
{\True $11,66\ \mathrm{kg\,m/s}$}
\loigiai{
Hai vectơ vuông góc nên $p=\sqrt{p_1^2+p_2^2}=\sqrt{10^2+6^2}$. Suy ra kết quả $11,66\ \mathrm{kg\,m/s}$.
}
\end{ex}

% ===== Câu 117 =====
% ID: AIQ_L10C5_FULL_0279
% Mức: VDC
\begin{ex}
% ID: AIQ_L10C5_FULL_0279
% Mức: VDC
Hai mảnh văng theo hai phương vuông góc có động lượng $11$ và $7\ \mathrm{kg\,m/s}$. Tính độ lớn tổng động lượng.
\choice
{\True $13,04\ \mathrm{kg\,m/s}$}
{$18\ \mathrm{kg\,m/s}$}
{$4\ \mathrm{kg\,m/s}$}
{$77\ \mathrm{kg\,m/s}$}
\loigiai{
Hai vectơ vuông góc nên $p=\sqrt{p_1^2+p_2^2}=\sqrt{11^2+7^2}$. Suy ra kết quả $13,04\ \mathrm{kg\,m/s}$.
}
\end{ex}

% ===== Câu 118 =====
% ID: AIQ_L10C5_FULL_0280
% Mức: VDC
\begin{ex}
% ID: AIQ_L10C5_FULL_0280
% Mức: VDC
Hai mảnh văng theo hai phương vuông góc có động lượng $12$ và $8\ \mathrm{kg\,m/s}$. Tính độ lớn tổng động lượng.
\choice
{$20\ \mathrm{kg\,m/s}$}
{\True $14,42\ \mathrm{kg\,m/s}$}
{$4\ \mathrm{kg\,m/s}$}
{$96\ \mathrm{kg\,m/s}$}
\loigiai{
Hai vectơ vuông góc nên $p=\sqrt{p_1^2+p_2^2}=\sqrt{12^2+8^2}$. Suy ra kết quả $14,42\ \mathrm{kg\,m/s}$.
}
\end{ex}

% ===== Câu 119 =====
% ID: AIQ_L10C5_FULL_0281
% Mức: NB
\begin{ex}
% ID: AIQ_L10C5_FULL_0281
% Mức: NB
Xét các phát biểu thuộc dạng “Bài toán động lượng hai chiều” ở tình huống số 1.
\choiceTF
{\True Động lượng hai chiều phải cộng vectơ.}
{\True Hai vectơ vuông góc dùng định lí Pythagore.}
{Có thể cộng đại số độ lớn cho mọi góc.}
{\True Cần xác định phương và chiều của hợp vectơ.}
\loigiai{
\begin{itemchoice}
\item (Đúng) Động lượng hai chiều phải cộng vectơ. (Vì): Phù hợp với định nghĩa hoặc hệ thức vật lí. b) Đúng: Phù hợp với định nghĩa hoặc hệ thức vật lí. c) Sai: Không phù hợp với định nghĩa hoặc hệ thức vật lí. d) Đúng: Phù hợp với định nghĩa hoặc hệ thức vật lí
\item (Đúng) Hai vectơ vuông góc dùng định lí Pythagore.
\item (Sai) Có thể cộng đại số độ lớn cho mọi góc.
\item (Đúng) Cần xác định phương và chiều của hợp vectơ.
\end{itemchoice}
}
\end{ex}

% ===== Câu 120 =====
% ID: AIQ_L10C5_FULL_0282
% Mức: TH
\begin{ex}
% ID: AIQ_L10C5_FULL_0282
% Mức: TH
Xét các phát biểu thuộc dạng “Bài toán động lượng hai chiều” ở tình huống số 2.
\choiceTF
{\True Động lượng hai chiều phải cộng vectơ.}
{\True Hai vectơ vuông góc dùng định lí Pythagore.}
{Có thể cộng đại số độ lớn cho mọi góc.}
{\True Cần xác định phương và chiều của hợp vectơ.}
\loigiai{
\begin{itemchoice}
\item (Đúng) Động lượng hai chiều phải cộng vectơ. (Vì): Phù hợp với định nghĩa hoặc hệ thức vật lí. b) Đúng: Phù hợp với định nghĩa hoặc hệ thức vật lí. c) Sai: Không phù hợp với định nghĩa hoặc hệ thức vật lí. d) Đúng: Phù hợp với định nghĩa hoặc hệ thức vật lí
\item (Đúng) Hai vectơ vuông góc dùng định lí Pythagore.
\item (Sai) Có thể cộng đại số độ lớn cho mọi góc.
\item (Đúng) Cần xác định phương và chiều của hợp vectơ.
\end{itemchoice}
}
\end{ex}

% ===== Câu 121 =====
% ID: AIQ_L10C5_FULL_0283
% Mức: TH
\begin{ex}
% ID: AIQ_L10C5_FULL_0283
% Mức: TH
Xét các phát biểu thuộc dạng “Bài toán động lượng hai chiều” ở tình huống số 3.
\choiceTF
{\True Động lượng hai chiều phải cộng vectơ.}
{\True Hai vectơ vuông góc dùng định lí Pythagore.}
{Có thể cộng đại số độ lớn cho mọi góc.}
{\True Cần xác định phương và chiều của hợp vectơ.}
\loigiai{
\begin{itemchoice}
\item (Đúng) Động lượng hai chiều phải cộng vectơ. (Vì): Phù hợp với định nghĩa hoặc hệ thức vật lí. b) Đúng: Phù hợp với định nghĩa hoặc hệ thức vật lí. c) Sai: Không phù hợp với định nghĩa hoặc hệ thức vật lí. d) Đúng: Phù hợp với định nghĩa hoặc hệ thức vật lí
\item (Đúng) Hai vectơ vuông góc dùng định lí Pythagore.
\item (Sai) Có thể cộng đại số độ lớn cho mọi góc.
\item (Đúng) Cần xác định phương và chiều của hợp vectơ.
\end{itemchoice}
}
\end{ex}

% ===== Câu 122 =====
% ID: AIQ_L10C5_FULL_0284
% Mức: VD
\begin{ex}
% ID: AIQ_L10C5_FULL_0284
% Mức: VD
Xét các phát biểu thuộc dạng “Bài toán động lượng hai chiều” ở tình huống số 4.
\choiceTF
{\True Động lượng hai chiều phải cộng vectơ.}
{\True Hai vectơ vuông góc dùng định lí Pythagore.}
{Có thể cộng đại số độ lớn cho mọi góc.}
{\True Cần xác định phương và chiều của hợp vectơ.}
\loigiai{
\begin{itemchoice}
\item (Đúng) Động lượng hai chiều phải cộng vectơ. (Vì): Phù hợp với định nghĩa hoặc hệ thức vật lí. b) Đúng: Phù hợp với định nghĩa hoặc hệ thức vật lí. c) Sai: Không phù hợp với định nghĩa hoặc hệ thức vật lí. d) Đúng: Phù hợp với định nghĩa hoặc hệ thức vật lí
\item (Đúng) Hai vectơ vuông góc dùng định lí Pythagore.
\item (Sai) Có thể cộng đại số độ lớn cho mọi góc.
\item (Đúng) Cần xác định phương và chiều của hợp vectơ.
\end{itemchoice}
}
\end{ex}

% ===== Câu 123 =====
% ID: AIQ_L10C5_FULL_0285
% Mức: VDC
\begin{ex}
% ID: AIQ_L10C5_FULL_0285
% Mức: VDC
Xét các phát biểu thuộc dạng “Bài toán động lượng hai chiều” ở tình huống số 5.
\choiceTF
{\True Động lượng hai chiều phải cộng vectơ.}
{\True Hai vectơ vuông góc dùng định lí Pythagore.}
{Có thể cộng đại số độ lớn cho mọi góc.}
{\True Cần xác định phương và chiều của hợp vectơ.}
\loigiai{
\begin{itemchoice}
\item (Đúng) Động lượng hai chiều phải cộng vectơ. (Vì): Phù hợp với định nghĩa hoặc hệ thức vật lí. b) Đúng: Phù hợp với định nghĩa hoặc hệ thức vật lí. c) Sai: Không phù hợp với định nghĩa hoặc hệ thức vật lí. d) Đúng: Phù hợp với định nghĩa hoặc hệ thức vật lí
\item (Đúng) Hai vectơ vuông góc dùng định lí Pythagore.
\item (Sai) Có thể cộng đại số độ lớn cho mọi góc.
\item (Đúng) Cần xác định phương và chiều của hợp vectơ.
\end{itemchoice}
}
\end{ex}

% ===== Câu 124 =====
% ID: AIQ_L10C5_FULL_0286
% Mức: TH
\begin{ex}
% ID: AIQ_L10C5_FULL_0286
% Mức: TH
Hai mảnh văng theo hai phương vuông góc có động lượng $13$ và $4\ \mathrm{kg\,m/s}$. Tính độ lớn tổng động lượng.
\shortans{13,6}
\loigiai{
Hai vectơ vuông góc nên $p=\sqrt{p_1^2+p_2^2}=\sqrt{13^2+4^2}$. Suy ra kết quả $13,6\ \mathrm{kg\,m/s}$. Đáp án trả lời ngắn không ghi đơn vị.
}
\end{ex}

% ===== Câu 125 =====
% ID: AIQ_L10C5_FULL_0287
% Mức: TH
\begin{ex}
% ID: AIQ_L10C5_FULL_0287
% Mức: TH
Hai mảnh văng theo hai phương vuông góc có động lượng $14$ và $5\ \mathrm{kg\,m/s}$. Tính độ lớn tổng động lượng.
\shortans{14,87}
\loigiai{
Hai vectơ vuông góc nên $p=\sqrt{p_1^2+p_2^2}=\sqrt{14^2+5^2}$. Suy ra kết quả $14,87\ \mathrm{kg\,m/s}$. Đáp án trả lời ngắn không ghi đơn vị.
}
\end{ex}

% ===== Câu 126 =====
% ID: AIQ_L10C5_FULL_0288
% Mức: VD
\begin{ex}
% ID: AIQ_L10C5_FULL_0288
% Mức: VD
Hai mảnh văng theo hai phương vuông góc có động lượng $15$ và $6\ \mathrm{kg\,m/s}$. Tính độ lớn tổng động lượng.
\shortans{16,16}
\loigiai{
Hai vectơ vuông góc nên $p=\sqrt{p_1^2+p_2^2}=\sqrt{15^2+6^2}$. Suy ra kết quả $16,16\ \mathrm{kg\,m/s}$. Đáp án trả lời ngắn không ghi đơn vị.
}
\end{ex}

% ===== Câu 127 =====
% ID: AIQ_L10C5_FULL_0289
% Mức: VD
\begin{ex}
% ID: AIQ_L10C5_FULL_0289
% Mức: VD
Hai mảnh văng theo hai phương vuông góc có động lượng $16$ và $7\ \mathrm{kg\,m/s}$. Tính độ lớn tổng động lượng.
\shortans{17,46}
\loigiai{
Hai vectơ vuông góc nên $p=\sqrt{p_1^2+p_2^2}=\sqrt{16^2+7^2}$. Suy ra kết quả $17,46\ \mathrm{kg\,m/s}$. Đáp án trả lời ngắn không ghi đơn vị.
}
\end{ex}

% ===== Câu 128 =====
% ID: AIQ_L10C5_FULL_0290
% Mức: VDC
\begin{ex}
% ID: AIQ_L10C5_FULL_0290
% Mức: VDC
Hai mảnh văng theo hai phương vuông góc có động lượng $17$ và $8\ \mathrm{kg\,m/s}$. Tính độ lớn tổng động lượng.
\shortans{18,79}
\loigiai{
Hai vectơ vuông góc nên $p=\sqrt{p_1^2+p_2^2}=\sqrt{17^2+8^2}$. Suy ra kết quả $18,79\ \mathrm{kg\,m/s}$. Đáp án trả lời ngắn không ghi đơn vị.
}
\end{ex}

% ===== Câu 129 =====
% ID: AIQ_L10C5_FULL_0291
% Mức: VDC
\begin{ex}
% ID: AIQ_L10C5_FULL_0291
% Mức: VDC
Hai mảnh văng theo hai phương vuông góc có động lượng $18$ và $4\ \mathrm{kg\,m/s}$. Tính độ lớn tổng động lượng.
\shortans{18,44}
\loigiai{
Hai vectơ vuông góc nên $p=\sqrt{p_1^2+p_2^2}=\sqrt{18^2+4^2}$. Suy ra kết quả $18,44\ \mathrm{kg\,m/s}$. Đáp án trả lời ngắn không ghi đơn vị.
}
\end{ex}

% ===== Câu 130 =====
% ID: AIQ_L10C5_FULL_0292
% Mức: TH
\begin{bt}
% ID: AIQ_L10C5_FULL_0292
% Mức: TH
Trình bày cách giải: Hai mảnh văng theo hai phương vuông góc có động lượng $19$ và $5\ \mathrm{kg\,m/s}$. Tính độ lớn tổng động lượng.
\loigiai{
Hai vectơ vuông góc nên $p=\sqrt{p_1^2+p_2^2}=\sqrt{19^2+5^2}$. Suy ra kết quả $19,65\ \mathrm{kg\,m/s}$.
}
\end{bt}

% ===== Câu 131 =====
% ID: AIQ_L10C5_FULL_0293
% Mức: TH
\begin{bt}
% ID: AIQ_L10C5_FULL_0293
% Mức: TH
Trình bày cách giải: Hai mảnh văng theo hai phương vuông góc có động lượng $20$ và $6\ \mathrm{kg\,m/s}$. Tính độ lớn tổng động lượng.
\loigiai{
Hai vectơ vuông góc nên $p=\sqrt{p_1^2+p_2^2}=\sqrt{20^2+6^2}$. Suy ra kết quả $20,88\ \mathrm{kg\,m/s}$.
}
\end{bt}

% ===== Câu 132 =====
% ID: AIQ_L10C5_FULL_0294
% Mức: VD
\begin{bt}
% ID: AIQ_L10C5_FULL_0294
% Mức: VD
Trình bày cách giải: Hai mảnh văng theo hai phương vuông góc có động lượng $21$ và $7\ \mathrm{kg\,m/s}$. Tính độ lớn tổng động lượng.
\loigiai{
Hai vectơ vuông góc nên $p=\sqrt{p_1^2+p_2^2}=\sqrt{21^2+7^2}$. Suy ra kết quả $22,14\ \mathrm{kg\,m/s}$.
}
\end{bt}

% ===== Câu 133 =====
% ID: AIQ_L10C5_FULL_0295
% Mức: VD
\begin{bt}
% ID: AIQ_L10C5_FULL_0295
% Mức: VD
Trình bày cách giải: Hai mảnh văng theo hai phương vuông góc có động lượng $22$ và $8\ \mathrm{kg\,m/s}$. Tính độ lớn tổng động lượng.
\loigiai{
Hai vectơ vuông góc nên $p=\sqrt{p_1^2+p_2^2}=\sqrt{22^2+8^2}$. Suy ra kết quả $23,41\ \mathrm{kg\,m/s}$.
}
\end{bt}

% ===== Câu 134 =====
% ID: AIQ_L10C5_FULL_0296
% Mức: VDC
\begin{bt}
% ID: AIQ_L10C5_FULL_0296
% Mức: VDC
Trình bày cách giải: Hai mảnh văng theo hai phương vuông góc có động lượng $23$ và $4\ \mathrm{kg\,m/s}$. Tính độ lớn tổng động lượng.
\loigiai{
Hai vectơ vuông góc nên $p=\sqrt{p_1^2+p_2^2}=\sqrt{23^2+4^2}$. Suy ra kết quả $23,35\ \mathrm{kg\,m/s}$.
}
\end{bt}

% ===== Câu 135 =====
% ID: AIQ_L10C5_FULL_0297
% Mức: VDC
\begin{bt}
% ID: AIQ_L10C5_FULL_0297
% Mức: VDC
Trình bày cách giải: Hai mảnh văng theo hai phương vuông góc có động lượng $24$ và $5\ \mathrm{kg\,m/s}$. Tính độ lớn tổng động lượng.
\loigiai{
Hai vectơ vuông góc nên $p=\sqrt{p_1^2+p_2^2}=\sqrt{24^2+5^2}$. Suy ra kết quả $24,52\ \mathrm{kg\,m/s}$.
}
\end{bt}

% ===== Câu 136 =====
% ID: AIQ_L10C5_FULL_0298
% Mức: NB
\dangbt{Bài toán tổng hợp bảo toàn động lượng và năng lượng}
\begin{ex}
% ID: AIQ_L10C5_FULL_0298
% Mức: NB
Hai vật cùng khối lượng, va chạm đàn hồi xuyên tâm; vật 1 ban đầu có tốc độ $3$ m/s, vật 2 đứng yên. Sau va chạm vật 1 dừng. Tốc độ vật 2 bằng bao nhiêu?
\choice
{\True $3\ \mathrm{m/s}$}
{$1,5\ \mathrm{m/s}$}
{$6\ \mathrm{m/s}$}
{$4\ \mathrm{m/s}$}
\loigiai{
Bảo toàn động lượng và động năng với hai khối lượng bằng nhau cho thấy hai vật trao đổi vận tốc. Suy ra kết quả $3\ \mathrm{m/s}$.
}
\end{ex}

% ===== Câu 137 =====
% ID: AIQ_L10C5_FULL_0299
% Mức: NB
\begin{ex}
% ID: AIQ_L10C5_FULL_0299
% Mức: NB
Hai vật cùng khối lượng, va chạm đàn hồi xuyên tâm; vật 1 ban đầu có tốc độ $4$ m/s, vật 2 đứng yên. Sau va chạm vật 1 dừng. Tốc độ vật 2 bằng bao nhiêu?
\choice
{$2\ \mathrm{m/s}$}
{\True $4\ \mathrm{m/s}$}
{$8\ \mathrm{m/s}$}
{$5\ \mathrm{m/s}$}
\loigiai{
Bảo toàn động lượng và động năng với hai khối lượng bằng nhau cho thấy hai vật trao đổi vận tốc. Suy ra kết quả $4\ \mathrm{m/s}$.
}
\end{ex}

% ===== Câu 138 =====
% ID: AIQ_L10C5_FULL_0300
% Mức: TH
\begin{ex}
% ID: AIQ_L10C5_FULL_0300
% Mức: TH
Hai vật cùng khối lượng, va chạm đàn hồi xuyên tâm; vật 1 ban đầu có tốc độ $5$ m/s, vật 2 đứng yên. Sau va chạm vật 1 dừng. Tốc độ vật 2 bằng bao nhiêu?
\choice
{$2,5\ \mathrm{m/s}$}
{$10\ \mathrm{m/s}$}
{\True $5\ \mathrm{m/s}$}
{$6\ \mathrm{m/s}$}
\loigiai{
Bảo toàn động lượng và động năng với hai khối lượng bằng nhau cho thấy hai vật trao đổi vận tốc. Suy ra kết quả $5\ \mathrm{m/s}$.
}
\end{ex}

% ===== Câu 139 =====
% ID: AIQ_L10C5_FULL_0301
% Mức: TH
\begin{ex}
% ID: AIQ_L10C5_FULL_0301
% Mức: TH
Hai vật cùng khối lượng, va chạm đàn hồi xuyên tâm; vật 1 ban đầu có tốc độ $6$ m/s, vật 2 đứng yên. Sau va chạm vật 1 dừng. Tốc độ vật 2 bằng bao nhiêu?
\choice
{$3\ \mathrm{m/s}$}
{$12\ \mathrm{m/s}$}
{$7\ \mathrm{m/s}$}
{\True $6\ \mathrm{m/s}$}
\loigiai{
Bảo toàn động lượng và động năng với hai khối lượng bằng nhau cho thấy hai vật trao đổi vận tốc. Suy ra kết quả $6\ \mathrm{m/s}$.
}
\end{ex}

% ===== Câu 140 =====
% ID: AIQ_L10C5_FULL_0302
% Mức: VD
\begin{ex}
% ID: AIQ_L10C5_FULL_0302
% Mức: VD
Hai vật cùng khối lượng, va chạm đàn hồi xuyên tâm; vật 1 ban đầu có tốc độ $7$ m/s, vật 2 đứng yên. Sau va chạm vật 1 dừng. Tốc độ vật 2 bằng bao nhiêu?
\choice
{\True $7\ \mathrm{m/s}$}
{$3,5\ \mathrm{m/s}$}
{$14\ \mathrm{m/s}$}
{$8\ \mathrm{m/s}$}
\loigiai{
Bảo toàn động lượng và động năng với hai khối lượng bằng nhau cho thấy hai vật trao đổi vận tốc. Suy ra kết quả $7\ \mathrm{m/s}$.
}
\end{ex}

% ===== Câu 141 =====
% ID: AIQ_L10C5_FULL_0303
% Mức: VD
\begin{ex}
% ID: AIQ_L10C5_FULL_0303
% Mức: VD
Hai vật cùng khối lượng, va chạm đàn hồi xuyên tâm; vật 1 ban đầu có tốc độ $8$ m/s, vật 2 đứng yên. Sau va chạm vật 1 dừng. Tốc độ vật 2 bằng bao nhiêu?
\choice
{$4\ \mathrm{m/s}$}
{\True $8\ \mathrm{m/s}$}
{$16\ \mathrm{m/s}$}
{$9\ \mathrm{m/s}$}
\loigiai{
Bảo toàn động lượng và động năng với hai khối lượng bằng nhau cho thấy hai vật trao đổi vận tốc. Suy ra kết quả $8\ \mathrm{m/s}$.
}
\end{ex}

% ===== Câu 142 =====
% ID: AIQ_L10C5_FULL_0304
% Mức: VD
\begin{ex}
% ID: AIQ_L10C5_FULL_0304
% Mức: VD
Hai vật cùng khối lượng, va chạm đàn hồi xuyên tâm; vật 1 ban đầu có tốc độ $9$ m/s, vật 2 đứng yên. Sau va chạm vật 1 dừng. Tốc độ vật 2 bằng bao nhiêu?
\choice
{$4,5\ \mathrm{m/s}$}
{$18\ \mathrm{m/s}$}
{\True $9\ \mathrm{m/s}$}
{$10\ \mathrm{m/s}$}
\loigiai{
Bảo toàn động lượng và động năng với hai khối lượng bằng nhau cho thấy hai vật trao đổi vận tốc. Suy ra kết quả $9\ \mathrm{m/s}$.
}
\end{ex}

% ===== Câu 143 =====
% ID: AIQ_L10C5_FULL_0305
% Mức: VDC
\begin{ex}
% ID: AIQ_L10C5_FULL_0305
% Mức: VDC
Hai vật cùng khối lượng, va chạm đàn hồi xuyên tâm; vật 1 ban đầu có tốc độ $10$ m/s, vật 2 đứng yên. Sau va chạm vật 1 dừng. Tốc độ vật 2 bằng bao nhiêu?
\choice
{$5\ \mathrm{m/s}$}
{$20\ \mathrm{m/s}$}
{$11\ \mathrm{m/s}$}
{\True $10\ \mathrm{m/s}$}
\loigiai{
Bảo toàn động lượng và động năng với hai khối lượng bằng nhau cho thấy hai vật trao đổi vận tốc. Suy ra kết quả $10\ \mathrm{m/s}$.
}
\end{ex}

% ===== Câu 144 =====
% ID: AIQ_L10C5_FULL_0306
% Mức: VDC
\begin{ex}
% ID: AIQ_L10C5_FULL_0306
% Mức: VDC
Hai vật cùng khối lượng, va chạm đàn hồi xuyên tâm; vật 1 ban đầu có tốc độ $11$ m/s, vật 2 đứng yên. Sau va chạm vật 1 dừng. Tốc độ vật 2 bằng bao nhiêu?
\choice
{\True $11\ \mathrm{m/s}$}
{$5,5\ \mathrm{m/s}$}
{$22\ \mathrm{m/s}$}
{$12\ \mathrm{m/s}$}
\loigiai{
Bảo toàn động lượng và động năng với hai khối lượng bằng nhau cho thấy hai vật trao đổi vận tốc. Suy ra kết quả $11\ \mathrm{m/s}$.
}
\end{ex}

% ===== Câu 145 =====
% ID: AIQ_L10C5_FULL_0307
% Mức: VDC
\begin{ex}
% ID: AIQ_L10C5_FULL_0307
% Mức: VDC
Hai vật cùng khối lượng, va chạm đàn hồi xuyên tâm; vật 1 ban đầu có tốc độ $12$ m/s, vật 2 đứng yên. Sau va chạm vật 1 dừng. Tốc độ vật 2 bằng bao nhiêu?
\choice
{$6\ \mathrm{m/s}$}
{\True $12\ \mathrm{m/s}$}
{$24\ \mathrm{m/s}$}
{$13\ \mathrm{m/s}$}
\loigiai{
Bảo toàn động lượng và động năng với hai khối lượng bằng nhau cho thấy hai vật trao đổi vận tốc. Suy ra kết quả $12\ \mathrm{m/s}$.
}
\end{ex}

% ===== Câu 146 =====
% ID: AIQ_L10C5_FULL_0308
% Mức: NB
\begin{ex}
% ID: AIQ_L10C5_FULL_0308
% Mức: NB
Xét các phát biểu thuộc dạng “Bài toán tổng hợp bảo toàn động lượng và năng lượng” ở tình huống số 1.
\choiceTF
{\True Va chạm đàn hồi bảo toàn động lượng và động năng.}
{Va chạm mềm bảo toàn cơ năng.}
{\True Khối lượng bằng nhau trong va chạm đàn hồi xuyên tâm có thể trao đổi vận tốc.}
{\True Cần kiểm tra đồng thời các định luật phù hợp.}
\loigiai{
\begin{itemchoice}
\item (Đúng) Va chạm đàn hồi bảo toàn động lượng và động năng. (Vì): Phù hợp với định nghĩa hoặc hệ thức vật lí. b) Sai: Không phù hợp với định nghĩa hoặc hệ thức vật lí. c) Đúng: Phù hợp với định nghĩa hoặc hệ thức vật lí. d) Đúng: Phù hợp với định nghĩa hoặc hệ thức vật lí
\item (Sai) Va chạm mềm bảo toàn cơ năng.
\item (Đúng) Khối lượng bằng nhau trong va chạm đàn hồi xuyên tâm có thể trao đổi vận tốc.
\item (Đúng) Cần kiểm tra đồng thời các định luật phù hợp.
\end{itemchoice}
}
\end{ex}

% ===== Câu 147 =====
% ID: AIQ_L10C5_FULL_0309
% Mức: TH
\begin{ex}
% ID: AIQ_L10C5_FULL_0309
% Mức: TH
Xét các phát biểu thuộc dạng “Bài toán tổng hợp bảo toàn động lượng và năng lượng” ở tình huống số 2.
\choiceTF
{\True Va chạm đàn hồi bảo toàn động lượng và động năng.}
{Va chạm mềm bảo toàn cơ năng.}
{\True Khối lượng bằng nhau trong va chạm đàn hồi xuyên tâm có thể trao đổi vận tốc.}
{\True Cần kiểm tra đồng thời các định luật phù hợp.}
\loigiai{
\begin{itemchoice}
\item (Đúng) Va chạm đàn hồi bảo toàn động lượng và động năng. (Vì): Phù hợp với định nghĩa hoặc hệ thức vật lí. b) Sai: Không phù hợp với định nghĩa hoặc hệ thức vật lí. c) Đúng: Phù hợp với định nghĩa hoặc hệ thức vật lí. d) Đúng: Phù hợp với định nghĩa hoặc hệ thức vật lí
\item (Sai) Va chạm mềm bảo toàn cơ năng.
\item (Đúng) Khối lượng bằng nhau trong va chạm đàn hồi xuyên tâm có thể trao đổi vận tốc.
\item (Đúng) Cần kiểm tra đồng thời các định luật phù hợp.
\end{itemchoice}
}
\end{ex}

% ===== Câu 148 =====
% ID: AIQ_L10C5_FULL_0310
% Mức: TH
\begin{ex}
% ID: AIQ_L10C5_FULL_0310
% Mức: TH
Xét các phát biểu thuộc dạng “Bài toán tổng hợp bảo toàn động lượng và năng lượng” ở tình huống số 3.
\choiceTF
{\True Va chạm đàn hồi bảo toàn động lượng và động năng.}
{Va chạm mềm bảo toàn cơ năng.}
{\True Khối lượng bằng nhau trong va chạm đàn hồi xuyên tâm có thể trao đổi vận tốc.}
{\True Cần kiểm tra đồng thời các định luật phù hợp.}
\loigiai{
\begin{itemchoice}
\item (Đúng) Va chạm đàn hồi bảo toàn động lượng và động năng. (Vì): Phù hợp với định nghĩa hoặc hệ thức vật lí. b) Sai: Không phù hợp với định nghĩa hoặc hệ thức vật lí. c) Đúng: Phù hợp với định nghĩa hoặc hệ thức vật lí. d) Đúng: Phù hợp với định nghĩa hoặc hệ thức vật lí
\item (Sai) Va chạm mềm bảo toàn cơ năng.
\item (Đúng) Khối lượng bằng nhau trong va chạm đàn hồi xuyên tâm có thể trao đổi vận tốc.
\item (Đúng) Cần kiểm tra đồng thời các định luật phù hợp.
\end{itemchoice}
}
\end{ex}

% ===== Câu 149 =====
% ID: AIQ_L10C5_FULL_0311
% Mức: VD
\begin{ex}
% ID: AIQ_L10C5_FULL_0311
% Mức: VD
Xét các phát biểu thuộc dạng “Bài toán tổng hợp bảo toàn động lượng và năng lượng” ở tình huống số 4.
\choiceTF
{\True Va chạm đàn hồi bảo toàn động lượng và động năng.}
{Va chạm mềm bảo toàn cơ năng.}
{\True Khối lượng bằng nhau trong va chạm đàn hồi xuyên tâm có thể trao đổi vận tốc.}
{\True Cần kiểm tra đồng thời các định luật phù hợp.}
\loigiai{
\begin{itemchoice}
\item (Đúng) Va chạm đàn hồi bảo toàn động lượng và động năng. (Vì): Phù hợp với định nghĩa hoặc hệ thức vật lí. b) Sai: Không phù hợp với định nghĩa hoặc hệ thức vật lí. c) Đúng: Phù hợp với định nghĩa hoặc hệ thức vật lí. d) Đúng: Phù hợp với định nghĩa hoặc hệ thức vật lí
\item (Sai) Va chạm mềm bảo toàn cơ năng.
\item (Đúng) Khối lượng bằng nhau trong va chạm đàn hồi xuyên tâm có thể trao đổi vận tốc.
\item (Đúng) Cần kiểm tra đồng thời các định luật phù hợp.
\end{itemchoice}
}
\end{ex}

% ===== Câu 150 =====
% ID: AIQ_L10C5_FULL_0312
% Mức: VDC
\begin{ex}
% ID: AIQ_L10C5_FULL_0312
% Mức: VDC
Xét các phát biểu thuộc dạng “Bài toán tổng hợp bảo toàn động lượng và năng lượng” ở tình huống số 5.
\choiceTF
{\True Va chạm đàn hồi bảo toàn động lượng và động năng.}
{Va chạm mềm bảo toàn cơ năng.}
{\True Khối lượng bằng nhau trong va chạm đàn hồi xuyên tâm có thể trao đổi vận tốc.}
{\True Cần kiểm tra đồng thời các định luật phù hợp.}
\loigiai{
\begin{itemchoice}
\item (Đúng) Va chạm đàn hồi bảo toàn động lượng và động năng. (Vì): Phù hợp với định nghĩa hoặc hệ thức vật lí. b) Sai: Không phù hợp với định nghĩa hoặc hệ thức vật lí. c) Đúng: Phù hợp với định nghĩa hoặc hệ thức vật lí. d) Đúng: Phù hợp với định nghĩa hoặc hệ thức vật lí
\item (Sai) Va chạm mềm bảo toàn cơ năng.
\item (Đúng) Khối lượng bằng nhau trong va chạm đàn hồi xuyên tâm có thể trao đổi vận tốc.
\item (Đúng) Cần kiểm tra đồng thời các định luật phù hợp.
\end{itemchoice}
}
\end{ex}

% ===== Câu 151 =====
% ID: AIQ_L10C5_FULL_0313
% Mức: TH
\begin{ex}
% ID: AIQ_L10C5_FULL_0313
% Mức: TH
Hai vật cùng khối lượng, va chạm đàn hồi xuyên tâm; vật 1 ban đầu có tốc độ $13$ m/s, vật 2 đứng yên. Sau va chạm vật 1 dừng. Tốc độ vật 2 bằng bao nhiêu?
\shortans{13}
\loigiai{
Bảo toàn động lượng và động năng với hai khối lượng bằng nhau cho thấy hai vật trao đổi vận tốc. Suy ra kết quả $13\ \mathrm{m/s}$. Đáp án trả lời ngắn không ghi đơn vị.
}
\end{ex}

% ===== Câu 152 =====
% ID: AIQ_L10C5_FULL_0314
% Mức: TH
\begin{ex}
% ID: AIQ_L10C5_FULL_0314
% Mức: TH
Hai vật cùng khối lượng, va chạm đàn hồi xuyên tâm; vật 1 ban đầu có tốc độ $14$ m/s, vật 2 đứng yên. Sau va chạm vật 1 dừng. Tốc độ vật 2 bằng bao nhiêu?
\shortans{14}
\loigiai{
Bảo toàn động lượng và động năng với hai khối lượng bằng nhau cho thấy hai vật trao đổi vận tốc. Suy ra kết quả $14\ \mathrm{m/s}$. Đáp án trả lời ngắn không ghi đơn vị.
}
\end{ex}

% ===== Câu 153 =====
% ID: AIQ_L10C5_FULL_0315
% Mức: VD
\begin{ex}
% ID: AIQ_L10C5_FULL_0315
% Mức: VD
Hai vật cùng khối lượng, va chạm đàn hồi xuyên tâm; vật 1 ban đầu có tốc độ $15$ m/s, vật 2 đứng yên. Sau va chạm vật 1 dừng. Tốc độ vật 2 bằng bao nhiêu?
\shortans{15}
\loigiai{
Bảo toàn động lượng và động năng với hai khối lượng bằng nhau cho thấy hai vật trao đổi vận tốc. Suy ra kết quả $15\ \mathrm{m/s}$. Đáp án trả lời ngắn không ghi đơn vị.
}
\end{ex}

% ===== Câu 154 =====
% ID: AIQ_L10C5_FULL_0316
% Mức: VD
\begin{ex}
% ID: AIQ_L10C5_FULL_0316
% Mức: VD
Hai vật cùng khối lượng, va chạm đàn hồi xuyên tâm; vật 1 ban đầu có tốc độ $16$ m/s, vật 2 đứng yên. Sau va chạm vật 1 dừng. Tốc độ vật 2 bằng bao nhiêu?
\shortans{16}
\loigiai{
Bảo toàn động lượng và động năng với hai khối lượng bằng nhau cho thấy hai vật trao đổi vận tốc. Suy ra kết quả $16\ \mathrm{m/s}$. Đáp án trả lời ngắn không ghi đơn vị.
}
\end{ex}

% ===== Câu 155 =====
% ID: AIQ_L10C5_FULL_0317
% Mức: VDC
\begin{ex}
% ID: AIQ_L10C5_FULL_0317
% Mức: VDC
Hai vật cùng khối lượng, va chạm đàn hồi xuyên tâm; vật 1 ban đầu có tốc độ $17$ m/s, vật 2 đứng yên. Sau va chạm vật 1 dừng. Tốc độ vật 2 bằng bao nhiêu?
\shortans{17}
\loigiai{
Bảo toàn động lượng và động năng với hai khối lượng bằng nhau cho thấy hai vật trao đổi vận tốc. Suy ra kết quả $17\ \mathrm{m/s}$. Đáp án trả lời ngắn không ghi đơn vị.
}
\end{ex}

% ===== Câu 156 =====
% ID: AIQ_L10C5_FULL_0318
% Mức: VDC
\begin{ex}
% ID: AIQ_L10C5_FULL_0318
% Mức: VDC
Hai vật cùng khối lượng, va chạm đàn hồi xuyên tâm; vật 1 ban đầu có tốc độ $18$ m/s, vật 2 đứng yên. Sau va chạm vật 1 dừng. Tốc độ vật 2 bằng bao nhiêu?
\shortans{18}
\loigiai{
Bảo toàn động lượng và động năng với hai khối lượng bằng nhau cho thấy hai vật trao đổi vận tốc. Suy ra kết quả $18\ \mathrm{m/s}$. Đáp án trả lời ngắn không ghi đơn vị.
}
\end{ex}

% ===== Câu 157 =====
% ID: AIQ_L10C5_FULL_0319
% Mức: TH
\begin{bt}
% ID: AIQ_L10C5_FULL_0319
% Mức: TH
Trình bày cách giải: Hai vật cùng khối lượng, va chạm đàn hồi xuyên tâm; vật 1 ban đầu có tốc độ $19$ m/s, vật 2 đứng yên. Sau va chạm vật 1 dừng. Tốc độ vật 2 bằng bao nhiêu?
\loigiai{
Bảo toàn động lượng và động năng với hai khối lượng bằng nhau cho thấy hai vật trao đổi vận tốc. Suy ra kết quả $19\ \mathrm{m/s}$.
}
\end{bt}

% ===== Câu 158 =====
% ID: AIQ_L10C5_FULL_0320
% Mức: TH
\begin{bt}
% ID: AIQ_L10C5_FULL_0320
% Mức: TH
Trình bày cách giải: Hai vật cùng khối lượng, va chạm đàn hồi xuyên tâm; vật 1 ban đầu có tốc độ $20$ m/s, vật 2 đứng yên. Sau va chạm vật 1 dừng. Tốc độ vật 2 bằng bao nhiêu?
\loigiai{
Bảo toàn động lượng và động năng với hai khối lượng bằng nhau cho thấy hai vật trao đổi vận tốc. Suy ra kết quả $20\ \mathrm{m/s}$.
}
\end{bt}

% ===== Câu 159 =====
% ID: AIQ_L10C5_FULL_0321
% Mức: VD
\begin{bt}
% ID: AIQ_L10C5_FULL_0321
% Mức: VD
Trình bày cách giải: Hai vật cùng khối lượng, va chạm đàn hồi xuyên tâm; vật 1 ban đầu có tốc độ $21$ m/s, vật 2 đứng yên. Sau va chạm vật 1 dừng. Tốc độ vật 2 bằng bao nhiêu?
\loigiai{
Bảo toàn động lượng và động năng với hai khối lượng bằng nhau cho thấy hai vật trao đổi vận tốc. Suy ra kết quả $21\ \mathrm{m/s}$.
}
\end{bt}

% ===== Câu 160 =====
% ID: AIQ_L10C5_FULL_0322
% Mức: VD
\begin{bt}
% ID: AIQ_L10C5_FULL_0322
% Mức: VD
Trình bày cách giải: Hai vật cùng khối lượng, va chạm đàn hồi xuyên tâm; vật 1 ban đầu có tốc độ $22$ m/s, vật 2 đứng yên. Sau va chạm vật 1 dừng. Tốc độ vật 2 bằng bao nhiêu?
\loigiai{
Bảo toàn động lượng và động năng với hai khối lượng bằng nhau cho thấy hai vật trao đổi vận tốc. Suy ra kết quả $22\ \mathrm{m/s}$.
}
\end{bt}

% ===== Câu 161 =====
% ID: AIQ_L10C5_FULL_0323
% Mức: VDC
\begin{bt}
% ID: AIQ_L10C5_FULL_0323
% Mức: VDC
Trình bày cách giải: Hai vật cùng khối lượng, va chạm đàn hồi xuyên tâm; vật 1 ban đầu có tốc độ $23$ m/s, vật 2 đứng yên. Sau va chạm vật 1 dừng. Tốc độ vật 2 bằng bao nhiêu?
\loigiai{
Bảo toàn động lượng và động năng với hai khối lượng bằng nhau cho thấy hai vật trao đổi vận tốc. Suy ra kết quả $23\ \mathrm{m/s}$.
}
\end{bt}

% ===== Câu 162 =====
% ID: AIQ_L10C5_FULL_0324
% Mức: VDC
\begin{bt}
% ID: AIQ_L10C5_FULL_0324
% Mức: VDC
Trình bày cách giải: Hai vật cùng khối lượng, va chạm đàn hồi xuyên tâm; vật 1 ban đầu có tốc độ $24$ m/s, vật 2 đứng yên. Sau va chạm vật 1 dừng. Tốc độ vật 2 bằng bao nhiêu?
\loigiai{
Bảo toàn động lượng và động năng với hai khối lượng bằng nhau cho thấy hai vật trao đổi vận tốc. Suy ra kết quả $24\ \mathrm{m/s}$.
}
\end{bt}
